\documentclass[11pt]{ctexart}
\usepackage[a4paper,margin=1in]{geometry}
\usepackage{graphicx}
\usepackage{xcolor}
\usepackage{hyperref}
\definecolor{refgray}{gray}{0.45}
\title{\textbf{市场最新观点汇总}\\[0.4em]{\large 油价回落与美元走强重塑亚洲央行政策，AI落地进入组织变革深水区，大宗商品主线从能源转向金属与电力}}
\author{KC桌面 自动生成}
\date{260629}
\begin{document}
\maketitle
\tableofcontents
\newpage
\section{一页摘要}
\begin{itemize}
  \item 亚洲新兴市场增长强劲，但油价下跌与美元走强正在系统性重塑各国央行的政策路径，科技出口仍为关键支撑。
  \item 大宗商品配置逻辑发生结构性转变，能源冲击后，铜、锂、铝及电力基础设施成为新的主线，供给约束长期存在。
  \item 美股融资市场发出二季度末最危险信号，AXW期货逼近历史极值，去杠杆风险上升，市场可能面临一次流动性事件。
  \item 存储器超级周期可能延续至2030年，由AI需求与长期协议驱动，行业正从周期性转向高利润、低波动模式。
  \item 人形机器人量产拐点已至，2026年中国市场出货量预期大幅上修，但真正的主线在于数据采集成本下降与代工模式成型。
  \item 财政风险冲击正制造滞胀陷阱，高债务经济体加息效果被削弱，央行宽松可能加剧通胀而非缓解衰退。
  \item AI价值分配出现结构性断裂，仅5\%的公司真正获得可量化回报，CEO亲自参与治理是区分赢家与输家的关键。
  \item 韩国军工与机器人板块下半年催化密集，西欧国家加速依赖韩国军工快速交付能力，机器人政策全球加码。
\end{itemize}
\section{宏观与利率：亚洲央行轮动与全球财政风险}
\textbf{油价下跌与美元走强正在重塑亚洲央行政策路径，同时BIS研究揭示高债务环境削弱加息效果，财政风险冲击可能引发滞胀。}

\begin{itemize}
  \item 亚洲新兴市场过去两个季度经济增速平均超过6\%（季调后年化），科技出口与生产持续爆发，但油价大幅回落与美元重新走强正在系统性改变各国央行的约束条件。
  \item JPM认为能源价格下行是下半年增长的超预期来源，通胀预期下降将提升居民购买力，菲律宾、印度、泰国、新加坡央行获得更大政策空间。
  \item 中国下半年增速预计放缓，但对亚洲的溢出效应有限，因亚洲对中国的出口依赖度已下降；但中国内需疲弱可能增加过剩产能外溢压力。
  \item 美元走强使部分央行压力上升，印尼盾承压可能迫使印尼央行继续加息，韩国央行面临科技热潮与美元走强的双重权衡。
  \item BIS工作论文发现，财政风险冲击会同时推高通胀和压低产出，形成滞胀格局；若央行用宽松货币政策对冲，结果将是通胀更高、衰退更深。
  \item BIS另一项研究显示，当公共债务占GDP比例从60\%上升到120\%时，同样幅度的加息对通胀的抑制效果显著减弱，而产出收缩幅度至少持平。
  \item 债务期限结构是关键变量：中间期限（9个月至8年）债务弱化货币政策效果，超短期与超长期债务反而放大效果。
  \item 欧元区加息对非欧元区欧洲国家有溢出效应，接收国4至8年期限债务占比越高，溢出效果越强。
\end{itemize}
\begin{figure}[htbp]
\centering
\includegraphics[width=0.88\linewidth]{figures/fig_001.jpg}
\caption{Figure 1 - Figure 1: EMAX tech IP and exports \%3m/3m, saar, both scales. IP thru Apr, exports incl. May forecast \#\# Oil Down, Dollar Up: A Rotation in Asian Central Banks With energy prices declining but the US Dollar strengthening, pressu}
\end{figure}
\begin{figure}[htbp]
\centering
\includegraphics[width=0.88\linewidth]{figures/fig_100.jpg}
\caption{Figure 1 - Figure 1: Estimated fiscal risk shocks for the U.S. in 2025 Notes: This figure shows the evolution of fiscal risk shocks in the US in 2025. The thick line represents the median target shock from the distribution of shocks recovere}
\end{figure}
\begin{figure}[htbp]
\centering
\includegraphics[width=0.88\linewidth]{figures/fig_102.jpg}
\caption{Figure 3 - Figure 3: Baseline 2SLS local-projection impulse response functions Notes: Impulse response functions to a 1-percentage-point fiscal-risk-driven increase in the 5-year government bond yield, estimated via 2SLS-LP as in equation (3}
\end{figure}
\begin{figure}[htbp]
\centering
\includegraphics[width=0.88\linewidth]{figures/fig_111.jpg}
\caption{Figure 1 - Figure 2: Response of key variables to a monetary policy shock, using asymmetric shocks. Note: The shaded area shows the 90\% confidence intervals. The responses are in \%; the size of the monetary policy shock is one standard devia}
\end{figure}
\begin{figure}[htbp]
\centering
\includegraphics[width=0.88\linewidth]{figures/fig_113.jpg}
\caption{Figure 2 - Figure 3: Euro area government debt residual maturity, by select maturity bins, relative to the size of GDP. Note: The line denotes the median, the box the interquartile range, and the whiskers the 25th (75th) percentiles minus (p}
\end{figure}
{\small\color{refgray}\textbf{References:} [R001] JPM：亚洲央行正经历一场“油跌美升”的轮动; [R011] 国际清算银行：财政风险冲击正在制造“滞胀陷阱”，而央行宽松只会让情况更糟; [R012] 国际清算银行：高债务正在悄悄削弱加息的通胀抑制力}

\section{AI与算力：存储器超级周期与AI落地组织变革}
\textbf{存储器超级周期由AI需求与长期协议驱动，可能延续至2030年；AI价值分配出现结构性断裂，CEO亲自参与治理是区分赢家与输家的关键。}

\begin{itemize}
  \item BofA基于美光财报判断，存储器超级周期可能延续至2027年甚至2030年，行业正从周期性转向类似台积电的高利润、低波动模式。
  \item 美光5月季度营收410亿美元（同比+346\%），毛利率85\%，8月季度指引500亿美元；DRAM和NAND价格均创历史新高。
  \item 长期协议（LTA）增加使行业利润更稳定，新建晶圆厂面临高成本、监管、水电供应等多重挑战，即便到2028年产能扩张也有限。
  \item HBM4已开始贡献收入，美光HBM4销售额达10亿美元；四大超大规模企业2026年资本支出合计约6500亿美元（同比+80\%）。
  \item 麦肯锡调研显示，88\%的组织在试验AI，但81\%未看到实质性利润增长；86\%的领导者认为组织未准备好将AI融入日常运营。
  \item CEO对AI治理的参与度是与企业EBIT影响相关度最高的因素之一；大企业更倾向让CEO亲自监督，中小企业多交给IT部门。
  \item BCG研究指出，仅5\%的“未来型”公司真正从AI获得可量化价值，其收入增速是落后者的1.7倍，股东回报是3.6倍。
  \item AI价值70\%集中在核心业务（销售、制造、供应链、定价），智能体AI在2025年已占AI价值的17\%，预计2028年达29\%。
\end{itemize}
\begin{figure}[htbp]
\centering
\includegraphics[width=0.88\linewidth]{figures/fig_027.jpg}
\caption{Exhibit 3 - Exhibit 3: Strong MoM growth in Jun (US\textbackslash{}\$25.5bn; +16\% MoM); more than 3x higher than 2025 year-average level Korea semis exports (mostly memory) – first 20 days of month (US\textbackslash{}\$bn) BofA GLOBAL RESEARCH Exhibit 5: Record-high sales}
\end{figure}
\begin{figure}[htbp]
\centering
\includegraphics[width=0.88\linewidth]{figures/fig_029.jpg}
\caption{Exhibit 4 - Exhibit 4: Up 188\% YoY in Jun; already five consecutive months of triple-digit growth Exhibit 6: Significant margin enhancement in May-Q (GM 85\%, OPM 81\%); mgmt. expects high-margin profile to continue Micron – Gross/OP margin tr}
\end{figure}
\begin{figure}[htbp]
\centering
\includegraphics[width=0.88\linewidth]{figures/fig_031.jpg}
\caption{Exhibit 7 - Exhibit 7: Robust price hike observed in May-Q (DRAM up low-60\% QoQ and NAND up mid-80\%) Exhibit 8: Record-high margins in May-Q (DRAM 81\%, NAND 78\%) Micron – DRAM and NAND OP margin trend \#\# Global memory forecasts Exhibit 9:}
\end{figure}
\begin{figure}[htbp]
\centering
\includegraphics[width=0.88\linewidth]{figures/fig_035.jpg}
\caption{Exhibit 29 - Exhibit 29: Substantial HBM content growth expected in 2025/2026-27, led by B300, Rubin, and Rubin Ultra NVIDIA GPUs HBM spec evolution; Rubin Ultra HBM content also up 3x vs 2026 1 \$\textasciicircum{}\{st\}\$ Rubin BofA GLOBAL RESEARCH Exhibit 30: 1}
\end{figure}
\begin{figure}[htbp]
\centering
\includegraphics[width=0.88\linewidth]{figures/fig_039.jpg}
\caption{Exhibit 34 - Exhibit 35: Amazon, Microsoft, Alphabet, Meta, and Oracle capex collectively to grow \textbackslash{}\textasciitilde{}80\% YoY in 2026 to \textbackslash{}\$650bn even after strong 65\% YoY growth in 2025; 2027/28 capex to hit \textbackslash{}\$800/950bn+ Combined capex of top 4 US tech companies}
\end{figure}
\begin{figure}[htbp]
\centering
\includegraphics[width=0.88\linewidth]{figures/fig_165.jpg}
\caption{Exhibit 1 - • Use Fit-for-Purpose Technology and Data 17 How to Accelerate AI Value Creation 19 Appendix • AI Definitions • Survey Methodology \# Are You Generating Value from AI? How much value is your company generating from your investments in AI? It’s a question more C}
\end{figure}
\begin{figure}[htbp]
\centering
\includegraphics[width=0.88\linewidth]{figures/fig_166.jpg}
\caption{Exhibit 2 - AI-driven value accrues over time, creating a compounding effect. (See Exhibit 2.) Future-built companies that moved early enjoy outsized benefits across financial and operational fronts, and this performance gap is widening. Future-built firms plan to spend 2}
\end{figure}
\begin{figure}[htbp]
\centering
\includegraphics[width=0.88\linewidth]{figures/fig_262.jpg}
\caption{Exhibit 1 - The survey findings also shed light on how organizations are structuring their AI deployment efforts. Some essential elements for deploying AI tend to be fully or partially centralized (Exhibit 1). For risk and compliance, as well as data governance, organizat}
\end{figure}
{\small\color{refgray}\textbf{References:} [R004] 美国银行：存储器超级周期可能延续到2030年; [R016] 波士顿咨询：95\%的公司正在错过AI真正的价值拐点; [R021] 麦肯锡：88\%企业在试AI，但81\%未见到利润; [R022] 麦肯锡：CEO不下场，AI永远只是成本}

\section{能源与大宗：从能源转向金属与电力}
\textbf{大宗商品配置逻辑发生结构性转变，能源冲击后，铜、锂、铝及电力基础设施成为新的主线，供给约束长期存在。}

\begin{itemize}
  \item GS指出，霍尔木兹冲突揭示能源市场韧性但未改变长期脆弱性；冲突后原油需求将因战略补库出现超过100万桶/日的反弹。
  \item 大宗商品配置价值未降低，但主角从传统油气转向电力、铜、锂和铝；电动车加速替代、可再生能源投资加码、AI竞赛推高电力需求。
  \item 铜的结构性缺口因政策加速，GS预计到2030年电网和电力基础设施将驱动超60\%的铜需求增长；美国钢铝关税预期使非美市场快速吃紧。
  \item Bernstein分析，斗山能源面临约2500亿美元的可参与核电项目总市场规模，其中约440亿美元可直接转化为设备合同。
  \item 斗山能源的订单节奏清晰：近看欧洲（波兰今年签设备合同），远看美国（2030年前后落地）；燃气轮机业务也在扩张。
  \item 斗山能源40多年核电设备制造经验构成核心壁垒，年订单有望从2025年的15万亿韩元增长到2030年的21万亿韩元以上。
  \item 钠离子电池在储能领域的成本拐点可能比预期更早到来，CATL推出全球首个公用事业级钠离子储能系统Tener Sodium。
  \item 钠离子电池在储能场景中能量密度不是核心矛盾，成本、安全性和循环寿命才是关键；CATL计划2026年9月开始交付。
\end{itemize}
\begin{figure}[htbp]
\centering
\includegraphics[width=0.88\linewidth]{figures/fig_002.jpg}
\caption{Exhibit 1 - Exhibit 1: The Strait of Hormuz disruption contributed to commodity returns outperforming other assets year to date, especially as a strong roll yield for oil combined with higher prices Importantly, the upside to energy prices d}
\end{figure}
\begin{figure}[htbp]
\centering
\includegraphics[width=0.88\linewidth]{figures/fig_003.jpg}
\caption{Exhibit 2 - Exhibit 3: Policy-risk-driven copper flows into the US drove the ex-US market into a deficit this year}
\end{figure}
\begin{figure}[htbp]
\centering
\includegraphics[width=0.88\linewidth]{figures/fig_007.jpg}
\caption{Exhibit 6 - Exhibit 6: Global EV sales have sharply accelerated during the Iran conflict... Exhibit 7: ...as have China solar panel exports Combined with the fact that power infrastructure can face bottlenecks, as has been the case in the}
\end{figure}
\begin{figure}[htbp]
\centering
\includegraphics[width=0.88\linewidth]{figures/fig_075.jpg}
\caption{EXHIBIT 1 - EXHIBIT 2: Continued growth in nuclear and gas equipment orders are expected to support company's valuation}
\end{figure}
\begin{figure}[htbp]
\centering
\includegraphics[width=0.88\linewidth]{figures/fig_080.jpg}
\caption{EXHIBIT 7 - EXHIBIT 8: We estimate around 35GW of nuclear project with USD250bn of TAM where Doosan could realistically participate}
\end{figure}
\begin{figure}[htbp]
\centering
\includegraphics[width=0.88\linewidth]{figures/fig_087.jpg}
\caption{EXHIBIT 14 - EXHIBIT 15: There has never been a better time to be a manufacturer of natural gas turbines. Gas turbine orders are set to increase to 100GW annually in the next few years which outpace current capacity of 60GW Global gas turbine or}
\end{figure}
{\small\color{refgray}\textbf{References:} [R002] GS：大宗商品真正的故事不在能源，而在金属与电力; [R005] Bernstein：钠离子电池的“奇点”不在车，在储能; [R009] Bernstein：斗山能源的2500亿美元核能订单，不只是“韩国制造”的故事}

\section{地产与消费：金融科技与游戏行业的结构性复苏}
\textbf{金融科技行业进入利润与规模必须同时兑现的新阶段，游戏行业由四大趋势碰撞驱动新一轮增长。}

\begin{itemize}
  \item BCG报告显示，2025年全球金融科技总收入突破5000亿美元（同比+22\%），增速是传统金融机构的四倍以上。
  \item 增长本质是盈利能力的修复：前85家上市金融科技公司平均EBITDA利润率提升4个百分点至20\%，74\%已实现盈利。
  \item 融资集中在后期轮次，E轮及以后融资额增长超210\%，种子轮收缩约10\%；投资者只愿为已证明规模化路径的公司买单。
  \item 交易与投资板块增长38\%，存款增长30\%，支付仍是最大收入贡献者；亚太增速最快（25\%）。
  \item BCG游戏报告预测，全球游戏收入将从2025年的约2600亿美元增长至2030年的约3300亿美元，复合年增长率约5\%。
  \item 云游戏进入主流分水岭：60\%的玩家尝试过云游戏，80\%体验正面；云游戏收入预计从15亿美元飙升至300亿美元。
  \item UGC创作者经济爆发，仅Fortnite和Roblox在2025年的创作者分成就超过15亿美元；40\%的玩家消费UGC比一年前更多。
  \item 应用商店开放将重塑移动游戏50\%全球收入的分发逻辑，开发者费用降低，自主分发能力大增。
\end{itemize}
\begin{figure}[htbp]
\centering
\includegraphics[width=0.88\linewidth]{figures/fig_141.jpg}
\caption{Exhibit 2 - \# Global Fintech Revenues Break Half a Trillion Dollars in 2025 Global fintech revenue by vertical (2021–2025, \textbackslash{}\$B) Global fintech revenue by region (2021–2025, \textbackslash{}\$B) Sources: S\&P Capital IQ; BCG FinTech Control Tower; BCG Banking and Insurance Revenue Pools; B}
\end{figure}
\begin{figure}[htbp]
\centering
\includegraphics[width=0.88\linewidth]{figures/fig_142.jpg}
\caption{EXHIBIT 3 - Fintechs account for \textbackslash{}\textasciitilde{}4\% of global financial services revenues TOTAL GLOBAL REVENUE, 2025 (\textbackslash{}\$) Fintech revenue growth outpaced incumbents across all verticals FINTECH VS. INCUMBENT REVENUE GROWTH YOY, 2024–2025 (\%) \$\textasciicircum{}\{1\}\$ Sources: S\&P Capital IQ; Pitchbook; B}
\end{figure}
\begin{figure}[htbp]
\centering
\includegraphics[width=0.88\linewidth]{figures/fig_146.jpg}
\caption{Exhibit 6 - Maturation is also evident in where capital is going. Equity funding rose 53\% to \textbackslash{}\$58 billion in 2025, but funding growth was not evenly distributed. (See Exhibit 6.) Trading and investment fintechs captured roughly one-third of all funding, up from about one-}
\end{figure}
\begin{figure}[htbp]
\centering
\includegraphics[width=0.88\linewidth]{figures/fig_177.jpg}
\caption{Exhibit 1 - 10 Platform Evolution: From Console Wars to Cloud Wars 13 UGC: Welcome to the New Creator Economy 16 App Stores Opening Up: A Revolution for Distribution 19 Improving Monetization: The New Math of Game Pricing 23 Growth Through Disruption 24 About the Global G}
\end{figure}
\begin{figure}[htbp]
\centering
\includegraphics[width=0.88\linewidth]{figures/fig_190.jpg}
\caption{EXHIBIT 9 - of gamers who tried it reported an overall positive experience However, gamers who use cloud gaming also play on other platforms HOW CLOUD GAMERS APPORTION THEIR GAMING TIME (\%) TIME DEVOTED TO CLOUD GAMING (\%) Sources: BCG Global Gaming Survey (N = 2,972); BC}
\end{figure}
\begin{figure}[htbp]
\centering
\includegraphics[width=0.88\linewidth]{figures/fig_191.jpg}
\caption{EXHIBIT 9 - Sources: BCG Global Gaming Survey (N = 2,972); BCG analysis. \#\# EXHIBIT 9 Cloud Gaming Is Ready for Liftoff as the Gaming Experience Improves CLOUD GAMING USERS (MILLIONS) Note: CAGR = compound annual growth rate. CLOUD GAMING MARKET (\textbackslash{}\$BILLIONS) \#\# Seizing th}
\end{figure}
\begin{figure}[htbp]
\centering
\includegraphics[width=0.88\linewidth]{figures/fig_192.jpg}
\caption{EXHIBIT 10 - \#\# EXHIBIT 10 \#\# Gamers Are Interacting With UGC, but Creators Are Still a Minority More than 40\% of gamers are consuming more UGC than they did a year ago... Q: Select how much you agree/disagree with the following statement: I consume more UGC now than I did}
\end{figure}
{\small\color{refgray}\textbf{References:} [R014] 波士顿咨询：金融科技重回增长，但赢家的门槛已经变了; [R017] 波士顿咨询：游戏行业的下一个增长引擎不是硬件，是平台碰撞}

\section{区域市场：韩国军工与机器人、人民币汇率信号}
\textbf{韩国军工与机器人板块下半年催化密集，西欧国家加速依赖韩国军工快速交付能力；人民币中间价模型透露政策信号切换。}

\begin{itemize}
  \item GS调研显示，韩国军工股二季度跑输大盘，但订单进展停滞、资金向存储芯片集中、全球军工同行走弱三个因素正在同时逆转。
  \item 过去一个月出现至少五个关键催化事件：瑞士转向韩国防空系统、西班牙敲定K9订单、LIG D\&A与莱茵金属签署MOU等。
  \item 西欧国家正将韩国视为传统军工供应链之外的“第二选择”，交付速度成为核心议价权。
  \item 韩国机器人板块受全球政策加码推动：韩国科技部拨款建训练中心、美国商务部长称机器人为“技术竞争的下一个战场”。
  \item NOM人民币中间价模型投影值从6.8166降至6.8077，确认短期贬值压力边际缓解；逆周期因子继续平滑波动。
  \item 模型投影值下降幅度大于市场预期，暗示外部压力正在减弱；央行面临“被动贬值”压力减轻，获得主动管理空间。
  \item 关键时间节点包括7月底政治局会议、10月国庆黄金周、11月APEC、12月中央经济工作会议及年底中美高层互动。
\end{itemize}
\begin{figure}[htbp]
\centering
\includegraphics[width=0.88\linewidth]{figures/fig_069.jpg}
\caption{Exhibit 1 - Exhibit 1 - Exhibit 4). Exhibit 1: Hanwha Aerospace: 12m fwd P/E Exhibit 2: Hyundai Rotem: 12m fwd P/E Exhibit 3: Hanwha Aerospace: GSe EPS}
\end{figure}
\begin{figure}[htbp]
\centering
\includegraphics[width=0.88\linewidth]{figures/fig_072.jpg}
\caption{Exhibit 3 - Exhibit 5: Recent robotics VC deal activity underscores interests in the robotics sector}
\end{figure}
\begin{figure}[htbp]
\centering
\includegraphics[width=0.88\linewidth]{figures/fig_074.jpg}
\caption{Exhibit 6 - Exhibit 6: How long a task AI can do (at \$50\textbackslash{}\%\$ success rate) is growing exponentially Each dot is a frontier AI model AI Task Completion Time Horizon \#\# Price Target Risks and Methodology - Hanwha Aerospace Our 12-month target p}
\end{figure}
{\small\color{refgray}\textbf{References:} [R006] GS：韩国军工和机器人，下半年最值得关注的两个板块; [R008] NOM：人民币中间价模型透露的政策信号，比点位更重要}

\section{风险偏好：融资市场压力与财政风险冲击}
\textbf{美股融资市场发出二季度末最危险信号，AXW期货逼近历史极值，去杠杆风险上升；BIS研究揭示财政风险冲击可能引发滞胀。}

\begin{itemize}
  \item MS报告指出，1个月期AXW期货合约飙升至+200个基点，逼近2024年12月历史最高水平，股票融资成本已近极值。
  \item 国债融资市场异常平静，隔夜GC回购报价仅3.75/3.72，与股票融资市场形成鲜明分化，表明压力集中在杠杆买家端。
  \item AXW期货近1个月标准差超过2个σ，历史上这种极端值出现后标普500往往出现均值回归。
  \item 融资压力持续到6月30日后可能意味着市场需要一次明显的去杠杆事件才能让融资成本正常化。
  \item BIS论文发现，财政风险冲击会同时推高通胀和压低产出，形成滞胀格局；若央行用宽松货币政策对冲，结果将是通胀更高、衰退更深。
  \item 财政风险冲击的识别策略巧妙：当主权债券收益率上升而安全公司债券收益率下降，两者走势背离正是财政风险冲击的独特指纹。
  \item 财政恶化推通胀的效果比财政改善压通胀的效果更强，长期国债收益率和股市对坏消息更敏感，这种不对称性值得关注。
\end{itemize}
\begin{figure}[htbp]
\centering
\includegraphics[width=0.88\linewidth]{figures/fig_013.jpg}
\caption{Exhibit 4 - Exhibit 3: July 2026 1-month AXW futures contract since February 23 Exhibit 4: Generic front 1-month AXW futures contract since December 2020 Several factors are likely also contributing to rich AXW levels, including issuance-r}
\end{figure}
\begin{figure}[htbp]
\centering
\includegraphics[width=0.88\linewidth]{figures/fig_015.jpg}
\caption{Exhibit 5 - Exhibit 5: Generic third 1-month AXW futures contract and the S\&P 500 since December 2020 Exhibit 6: Generic third 1-month AXW futures contract and the S\&P 500 since June 2022, rolling 252-day z-score We think risk of a tumultu}
\end{figure}
\begin{figure}[htbp]
\centering
\includegraphics[width=0.88\linewidth]{figures/fig_100.jpg}
\caption{Figure 1 - Figure 1: Estimated fiscal risk shocks for the U.S. in 2025 Notes: This figure shows the evolution of fiscal risk shocks in the US in 2025. The thick line represents the median target shock from the distribution of shocks recovere}
\end{figure}
\begin{figure}[htbp]
\centering
\includegraphics[width=0.88\linewidth]{figures/fig_102.jpg}
\caption{Figure 3 - Figure 3: Baseline 2SLS local-projection impulse response functions Notes: Impulse response functions to a 1-percentage-point fiscal-risk-driven increase in the 5-year government bond yield, estimated via 2SLS-LP as in equation (3}
\end{figure}
\begin{figure}[htbp]
\centering
\includegraphics[width=0.88\linewidth]{figures/fig_104.jpg}
\caption{Figure 5 - Figure 5: State-dependent effects of fiscal risk shocks: high vs. low CDS environment (p90 vs. p10) Notes: Impulse response functions from equation (11), scaled by \$(z\_\{90\}-z\_\{10\})\$ to represent the difference in macro-financial t}
\end{figure}
{\small\color{refgray}\textbf{References:} [R003] 摩根斯坦利：MS：融资市场正发出二季度末最危险的信号; [R011] 国际清算银行：财政风险冲击正在制造“滞胀陷阱”，而央行宽松只会让情况更糟}

\section{人形机器人：量产拐点与反射层架构}
\textbf{人形机器人量产拐点已至，2026年中国市场出货量预期大幅上修，但真正的主线在于数据采集成本下降与代工模式成型。}

\begin{itemize}
  \item NOM预计2026年中国市场人形机器人出货量达4万至5万台，较此前预期显著上调，消费端而非工业端成为主要驱动力。
  \item 特斯拉Optimus产能规划上调：弗里蒙特产线2027年目标从5万台调至7万台，2028年奥斯汀第二条产线再增7万台。
  \item 下游结构中消费端和表演/娱乐端各占约30\%，政府采购占20\%，教育占15\%，商业和工业应用仅占3\%至5\%。
  \item 数据采集范式正在切换：行业从“真人机器人采集”转向“非机器人采集”，成本仅为真实采集的20\%，速度大幅提升。
  \item NOM提出物理AI的瓶颈不是大模型推理能力，而是能否构建低延迟的“反射层”（40毫秒响应比300毫秒深度思考更关键）。
  \item NeuralAxis架构将系统拆分为推理层（300ms）、协调层（小脑）、反射层（40ms），反射处理器分布到关节、手部和足部。
  \item 宇树科技WVLA2.0模型是反射层架构的首次商业化落地，融合WMA预测模型与VLA端到端动作生成。
  \item 代工趋势明确，硬件供应链重资产、回收期长，多数中小厂商选择外协，代工链条中不可替代的位置成为竞争焦点。
\end{itemize}
\begin{figure}[htbp]
\centering
\includegraphics[width=0.88\linewidth]{figures/fig_235.jpg}
\caption{Exhibit 1 - By contrast, a decline in the importance of software development and other highly specialized technical execution skills reflects that AI systems increasingly take over those tasks. At the same time, skills necessary to guide, interpret, and apply AI outputs a}
\end{figure}
\begin{figure}[htbp]
\centering
\includegraphics[width=0.88\linewidth]{figures/fig_236.jpg}
\caption{Exhibit 2 - \#\# Workforce planning remains predominantly short term and operational On average, 62 percent of HR professionals indicate that their companies conduct organization-wide workforce planning, while 34 percent apply it to at least part of their workforce. However}
\end{figure}
\begin{figure}[htbp]
\centering
\includegraphics[width=0.88\linewidth]{figures/fig_237.jpg}
\caption{Exhibit 2 - \#\# Exhibit 2 Engagement in workforce planning, \$\textasciicircum{}\{1\}\$ \% of respondents How far in advance companies typically forecast workforce needs, \$\textasciicircum{}\{4\}\$ \% of respondents Note: Figures may not sum to 100\%, because of rounding. \$\textasciicircum{}\{1\}\$ Question: Does your company carry out}
\end{figure}
{\small\color{refgray}\textbf{References:} [R007] NOM：人形机器人量产拐点已至，但真正的主线不在机器人本身; [R010] NOM：人形机器人的“脊髓”比“大脑”更关键}

\section{企业AI转型：CEO亲自下场与组织重塑}
\textbf{AI价值分配出现结构性断裂，CEO亲自参与治理是区分赢家与输家的关键；企业需推动技术与组织的双重转型。}

\begin{itemize}
  \item BCG报告显示，72\%的CEO直接负责公司AI决策，但仅15\%获得有意义的回报；先行者每周至少花8小时亲自使用AI。
  \item 先行者CEO用AI扮演四种角色：随需导师、魔鬼代言人、战略编辑、反思工具，核心是放大判断力而非节省时间。
  \item 麦肯锡调研发现，88\%的组织在试验AI，但81\%未看到利润增长；86\%的领导者认为组织未准备好将AI融入日常运营。
  \item CEO对AI治理的参与度是与企业EBIT影响相关度最高的因素之一；大企业更倾向让CEO亲自监督AI治理。
  \item 重新设计工作流程是对AI利润影响最大的组织变革动作，但仅21\%的企业真正动手改了流程。
  \item BCG指出，5\%的“未来型”公司正在构筑AI价值护城河，其收入增速是落后者的1.7倍，股东回报是3.6倍。
  \item 未来型公司计划将IT预算中64\%投入AI，而追赶者不足一半；AI价值70\%集中在核心业务功能。
  \item 智能体AI在2025年已占AI价值的17\%，预计2028年达29\%；智能体不是独立项目，需嵌入整个工作流。
\end{itemize}
\begin{figure}[htbp]
\centering
\includegraphics[width=0.88\linewidth]{figures/fig_165.jpg}
\caption{Exhibit 1 - • Use Fit-for-Purpose Technology and Data 17 How to Accelerate AI Value Creation 19 Appendix • AI Definitions • Survey Methodology \# Are You Generating Value from AI? How much value is your company generating from your investments in AI? It’s a question more C}
\end{figure}
\begin{figure}[htbp]
\centering
\includegraphics[width=0.88\linewidth]{figures/fig_166.jpg}
\caption{Exhibit 2 - AI-driven value accrues over time, creating a compounding effect. (See Exhibit 2.) Future-built companies that moved early enjoy outsized benefits across financial and operational fronts, and this performance gap is widening. Future-built firms plan to spend 2}
\end{figure}
\begin{figure}[htbp]
\centering
\includegraphics[width=0.88\linewidth]{figures/fig_168.jpg}
\caption{EXHIBIT 5 - Agents do not operate without humans. Their performance depends on strong human orchestration in redesigned roles. The best companies reconfigure workflows to combine autonomous agents with human oversight, maximizing value and adoption. \#\# EXHIBIT 5 \#\# Value }
\end{figure}
\begin{figure}[htbp]
\centering
\includegraphics[width=0.88\linewidth]{figures/fig_262.jpg}
\caption{Exhibit 1 - The survey findings also shed light on how organizations are structuring their AI deployment efforts. Some essential elements for deploying AI tend to be fully or partially centralized (Exhibit 1). For risk and compliance, as well as data governance, organizat}
\end{figure}
\begin{figure}[htbp]
\centering
\includegraphics[width=0.88\linewidth]{figures/fig_263.jpg}
\caption{Exhibit 2 - Organizations have employees overseeing the quality of gen AI outputs, though the extent of that oversight varies widely. Twenty-seven percent of respondents whose organizations use gen AI say that employees review all content created by gen AI before it is us}
\end{figure}
\begin{figure}[htbp]
\centering
\includegraphics[width=0.88\linewidth]{figures/fig_268.jpg}
\caption{Exhibit 9 - Organizations are also using AI in more business functions than in the previous State of AI survey. For the first time, most survey respondents report the use of AI in more than one business function (Exhibit 9). Responses show organizations using AI in an ave}
\end{figure}
{\small\color{refgray}\textbf{References:} [R013] 波士顿咨询：CEO的AI红利，不在效率，在判断; [R016] 波士顿咨询：95\%的公司正在错过AI真正的价值拐点; [R021] 麦肯锡：88\%企业在试AI，但81\%未见到利润; [R022] 麦肯锡：CEO不下场，AI永远只是成本}

\section{美国竞争力与劳动力市场转型}
\textbf{美国经济竞争力账面强劲但结构性裂缝清晰可见，HR职能正站在被IT吞噬与重塑组织的分叉口。}

\begin{itemize}
  \item 麦肯锡报告显示，美国以全球4\%的人口创造26\%的全球GDP，贡献超50\%的全球市值，但收入不平等、教育下滑、基建评级C级构成隐忧。
  \item 美国竞争力的三维定义中，全球领先企业与创新技术领导力突出，但经济机会维度出现明显裂痕。
  \item 麦肯锡HR Monitor报告指出，仅11\%的企业在做三年以上的战略性人力规划，而技能缺口正在加速扩大。
  \item 23\%的员工当前技能就不够用，22\%担心5年内跟不上；HR对“未来技能”的判断存在错位，软件工程需求在降，数字素养与推理能力在上升。
  \item HR职能正站在分叉路口：一条路是维持现状导致核心职责被IT或数字化部门蚕食，另一条路是主动重塑成为AI时代组织转型的架构师。
  \item 员工离职率下降2个百分点，但薪酬（52\%）、工作生活平衡（46\%）、工作稳定（45\%）是留任的核心因素。
  \item AI正在重塑岗位技能需求，HR需要从“管人”转向思考“人和智能体如何协作”的新命题。
\end{itemize}
\begin{figure}[htbp]
\centering
\includegraphics[width=0.88\linewidth]{figures/fig_207.jpg}
\caption{Exhibit 1 - Through every chapter of the past 250 years, the United States has harnessed these foundations, not in a fixed economic model but through flexible institutions that have made the next adaptation possible. And it has done so collectively. “We the people”—farmer}
\end{figure}
\begin{figure}[htbp]
\centering
\includegraphics[width=0.88\linewidth]{figures/fig_210.jpg}
\caption{Exhibit 3 - American companies make up more than half of the top 100 firms globally by market capitalization and revenue (Exhibit 3). From start-ups to large corporations, they attract an outsize share of capital from global markets. US firms hold more than half of global}
\end{figure}
\begin{figure}[htbp]
\centering
\includegraphics[width=0.88\linewidth]{figures/fig_214.jpg}
\caption{Exhibit 7 - This gap in income levels has grown over the past 50 years. Although all income segments have seen real growth, market incomes (wages and asset flows) have grown the most for the top two quintiles. For the bottom 60 percent of the population, more income growt}
\end{figure}
\begin{figure}[htbp]
\centering
\includegraphics[width=0.88\linewidth]{figures/fig_235.jpg}
\caption{Exhibit 1 - By contrast, a decline in the importance of software development and other highly specialized technical execution skills reflects that AI systems increasingly take over those tasks. At the same time, skills necessary to guide, interpret, and apply AI outputs a}
\end{figure}
\begin{figure}[htbp]
\centering
\includegraphics[width=0.88\linewidth]{figures/fig_236.jpg}
\caption{Exhibit 2 - \#\# Workforce planning remains predominantly short term and operational On average, 62 percent of HR professionals indicate that their companies conduct organization-wide workforce planning, while 34 percent apply it to at least part of their workforce. However}
\end{figure}
\begin{figure}[htbp]
\centering
\includegraphics[width=0.88\linewidth]{figures/fig_243.jpg}
\caption{Exhibit 9 - — Relationship with colleagues (27 percent, down six percentage points). A sense of belonging and positive team dynamics continue to matter, even though this category ranks below economic and structural factors. \#\# Exhibit 9 Employees now prioritize remunerati}
\end{figure}
{\small\color{refgray}\textbf{References:} [R018] 麦肯锡：美国250年竞争力未断，但真正考验的是“下一章”能否续写; [R019] 麦肯锡：HR正站在“被IT吞噬”与“重塑组织”的分叉口}

\section{金融科技与保险：盈利修复与市场分化}
\textbf{金融科技行业进入利润与规模必须同时兑现的新阶段，美国P\&C保险市场硬市场红利消退，头部与尾部公司差距拉大。}

\begin{itemize}
  \item BCG报告显示，2025年全球金融科技总收入突破5000亿美元（同比+22\%），增速是传统金融机构的四倍以上。
  \item 前85家上市金融科技公司平均EBITDA利润率提升4个百分点至20\%，74\%已实现盈利；融资集中在后期轮次。
  \item 交易与投资板块增长38\%，存款增长30\%，支付仍是最大收入贡献者；亚太增速最快（25\%）。
  \item 美国P\&C保险市场2021至2025年实现约18\%的年化总股东回报率，但2025年势头明显放缓。
  \item 头部P\&C公司承保业务对RoTE贡献约11个百分点，而尾部公司承保为负值；承保能力是分水岭，投资收入只是止痛药。
  \item 社会通胀（如“天价判决”）、强对流风暴损失、AI带来的新风险正在重塑风险组合。
  \item 领先保险公司已规模化部署AI于承保、理赔和定价，获得运营效率和更精细的风险洞察。
\end{itemize}
\begin{figure}[htbp]
\centering
\includegraphics[width=0.88\linewidth]{figures/fig_141.jpg}
\caption{Exhibit 2 - \# Global Fintech Revenues Break Half a Trillion Dollars in 2025 Global fintech revenue by vertical (2021–2025, \textbackslash{}\$B) Global fintech revenue by region (2021–2025, \textbackslash{}\$B) Sources: S\&P Capital IQ; BCG FinTech Control Tower; BCG Banking and Insurance Revenue Pools; B}
\end{figure}
\begin{figure}[htbp]
\centering
\includegraphics[width=0.88\linewidth]{figures/fig_142.jpg}
\caption{EXHIBIT 3 - Fintechs account for \textbackslash{}\textasciitilde{}4\% of global financial services revenues TOTAL GLOBAL REVENUE, 2025 (\textbackslash{}\$) Fintech revenue growth outpaced incumbents across all verticals FINTECH VS. INCUMBENT REVENUE GROWTH YOY, 2024–2025 (\%) \$\textasciicircum{}\{1\}\$ Sources: S\&P Capital IQ; Pitchbook; B}
\end{figure}
\begin{figure}[htbp]
\centering
\includegraphics[width=0.88\linewidth]{figures/fig_146.jpg}
\caption{Exhibit 6 - Maturation is also evident in where capital is going. Equity funding rose 53\% to \textbackslash{}\$58 billion in 2025, but funding growth was not evenly distributed. (See Exhibit 6.) Trading and investment fintechs captured roughly one-third of all funding, up from about one-}
\end{figure}
\begin{figure}[htbp]
\centering
\includegraphics[width=0.88\linewidth]{figures/fig_160.jpg}
\caption{Exhibit 1 - The primary engine of value creation over the 2021–2025 period remains growth in tangible book value (TBV), the result of strong underwriting income. (See Exhibit 1.) Cash flow contribution (dividends and buybacks) remains the second major lever, while multipl}
\end{figure}
\begin{figure}[htbp]
\centering
\includegraphics[width=0.88\linewidth]{figures/fig_161.jpg}
\caption{EXHIBIT 2 - Top quartile \#\# EXHIBIT 2 Underwriting Results Drove Performance of Top-Quartile Companies CONTRIBUTION TO RETURN ON TANGIBLE EQUITY, 2021–2025 (PP) Sources: S\&P Capital IQ; BCG ValueScience Center; BCG analysis. Note: PP = percentage points; RoTE = return on }
\end{figure}
\begin{figure}[htbp]
\centering
\includegraphics[width=0.88\linewidth]{figures/fig_163.jpg}
\caption{Exhibit 4 - The subsegments of US P\&C performed similarly from 2021 through 2025, with TSRs of 17\% and 18\% for personal and commercial lines, respectively. Behind these figures, however, are signs of diverging markets. (See Exhibit 4.) Commercial lines benefited from a pr}
\end{figure}
{\small\color{refgray}\textbf{References:} [R014] 波士顿咨询：金融科技重回增长，但赢家的门槛已经变了; [R015] 波士顿咨询：美国P\&C保险的“好日子”结束了，但赢家还没定}

\section{结语}
综合来看，当前市场主线清晰分化：亚洲央行在油价与美元轮动中寻找新平衡，AI落地进入组织变革深水区，大宗商品主线从能源转向金属与电力，而融资市场与财政风险则构成尾部风险。投资者需关注这些结构性变化而非短期波动。
\newpage
\section{报告覆盖清单}
\begin{itemize}
  \item [R001] 宏观与利率：亚洲央行轮动与全球财政风险 - JPM：亚洲央行正经历一场“油跌美升”的轮动
  \item [R002] 能源与大宗：从能源转向金属与电力 - GS：大宗商品真正的故事不在能源，而在金属与电力
  \item [R003] 风险偏好：融资市场压力与财政风险冲击 - 摩根斯坦利：MS：融资市场正发出二季度末最危险的信号
  \item [R004] AI与算力：存储器超级周期与AI落地组织变革 - 美国银行：存储器超级周期可能延续到2030年
  \item [R005] 能源与大宗：从能源转向金属与电力 - Bernstein：钠离子电池的“奇点”不在车，在储能
  \item [R006] 区域市场：韩国军工与机器人、人民币汇率信号 - GS：韩国军工和机器人，下半年最值得关注的两个板块
  \item [R007] 人形机器人：量产拐点与反射层架构 - NOM：人形机器人量产拐点已至，但真正的主线不在机器人本身
  \item [R008] 区域市场：韩国军工与机器人、人民币汇率信号 - NOM：人民币中间价模型透露的政策信号，比点位更重要
  \item [R009] 能源与大宗：从能源转向金属与电力 - Bernstein：斗山能源的2500亿美元核能订单，不只是“韩国制造”的故事
  \item [R010] 人形机器人：量产拐点与反射层架构 - NOM：人形机器人的“脊髓”比“大脑”更关键
  \item [R011] 宏观与利率：亚洲央行轮动与全球财政风险 / 风险偏好：融资市场压力与财政风险冲击 - 国际清算银行：财政风险冲击正在制造“滞胀陷阱”，而央行宽松只会让情况更糟
  \item [R012] 宏观与利率：亚洲央行轮动与全球财政风险 - 国际清算银行：高债务正在悄悄削弱加息的通胀抑制力
  \item [R013] 企业AI转型：CEO亲自下场与组织重塑 - 波士顿咨询：CEO的AI红利，不在效率，在判断
  \item [R014] 地产与消费：金融科技与游戏行业的结构性复苏 / 金融科技与保险：盈利修复与市场分化 - 波士顿咨询：金融科技重回增长，但赢家的门槛已经变了
  \item [R015] 金融科技与保险：盈利修复与市场分化 - 波士顿咨询：美国P\&C保险的“好日子”结束了，但赢家还没定
  \item [R016] AI与算力：存储器超级周期与AI落地组织变革 / 企业AI转型：CEO亲自下场与组织重塑 - 波士顿咨询：95\%的公司正在错过AI真正的价值拐点
  \item [R017] 地产与消费：金融科技与游戏行业的结构性复苏 - 波士顿咨询：游戏行业的下一个增长引擎不是硬件，是平台碰撞
  \item [R018] 美国竞争力与劳动力市场转型 - 麦肯锡：美国250年竞争力未断，但真正考验的是“下一章”能否续写
  \item [R019] 美国竞争力与劳动力市场转型 - 麦肯锡：HR正站在“被IT吞噬”与“重塑组织”的分叉口
  \item [R020] 未被主题摘要引用 - 麦肯锡：2025年，企业最大的风险不是AI，而是把学习当成福利
  \item [R021] AI与算力：存储器超级周期与AI落地组织变革 / 企业AI转型：CEO亲自下场与组织重塑 - 麦肯锡：88\%企业在试AI，但81\%未见到利润
  \item [R022] AI与算力：存储器超级周期与AI落地组织变革 / 企业AI转型：CEO亲自下场与组织重塑 - 麦肯锡：CEO不下场，AI永远只是成本
\end{itemize}
\section{逐篇报告摘录}
\subsection{[R001] JPM：亚洲央行正经历一场“油跌美升”的轮动}
{\small\color{refgray}归类：宏观与利率：亚洲央行轮动与全球财政风险}

过去两个月，亚洲新兴市场（EMAX）经济增速平均超过6\%（季调后年化），几乎是潜在增长率的两倍。支撑这一强劲表现的是科技出口与生产的持续爆发。但这份JPM最新发布的《全球数据观察：亚洲》报告提出了一个更值得关注的结构性判断：随着中东冲突缓和、原油价格大幅回落，叠加美元重新走强，亚洲各经济体央行面临的约束条件正在发生系统性轮动。对于投资者而言，真正重要的不是当前的增长有多强，而是这种“油跌美升”的格局将如何重塑各国央行的政策路径，并由此改变不同市场的资产定价逻辑。

这份报告的核心主张是：能源价格下行正在为亚洲经济下半年的增长带来上行风险，但美元走强同时也在制造新的分化——一些央行将因此获得喘息空间，另一些则可能被迫进一步收紧。理解这场轮动，是判断下半年亚洲资产配置方向的关键前提。

报告明确指出，中东冲突的缓和与原油价格的急剧下跌，显著缓解了亚洲地区此前面临的贸易条件恶化冲击。在此之前，亚洲经济体对中东能源的依赖使其在地缘冲突中尤为脆弱。但现实比预期更有韧性——各国通过积极从其他渠道采购能源、消耗储备以及转向替代品，避免了最坏的结果。

这意味着，尽管能源逆风依然存在，二季度增长仍接近趋势水平。而展望下半年，报告认为上行风险正在积聚。科技生产继续保持强劲势头——虽然台湾5月工业产出显示科技动能较此前“滚烫”的速度有所降温，但仍在以30\%的年化季环比增长运行；韩国20天出口数据同样指向强劲的出口动能，尽管近期科技涨幅更多由价格驱动。

KC评论： 市场对亚洲增长的担忧主要集中在中国的放缓上，但JPM认为能源价格的下跌正在成为一个被低估的正面因素。通胀预期下降将提升家庭购买力，减少部分经济体（如菲律宾）所需的货币紧缩幅度，并强化其他经济体（印度、泰国、新加坡）央行无需加息的判断。这意味着，对于那些能源进口依赖度高且国内需求对利率敏感的经济体，下半年可能迎来一个“成本下降+金融条件放松”的双重利好窗口。完整报告中关于各国通胀路径的具体假设和敏感性分析，值得仔细推敲。

报告承认，中国下半年增长的急剧减速——JPM预测二季度GDP季调后年化增速将放缓至3.3\%——将为区域增长带来逆风。但关键在于，这种冲击正在被结构性因素分散化。

近年来，该地区对中国的出口敞口已经降低。报告特别指出，只有印尼和菲律宾在中国近几个季度超趋势增长期间，对华出口出现了有意义的上升。因此，中国增速放缓的影响可能只集中在少数几个经济体。更值得警惕的是另一种传导机制：如果中国国内需求持续疲弱，将增加其向区域乃至全球出口过剩产能的激励。这意味着，对于东南亚出口导向型经济体而言，来自中国的竞争压力可能比需求萎缩更为持久。

KC评论： 这里有一个容易被忽略的二阶效应：中国产能过剩的溢出，对于依赖制造业出口的亚洲经济体既是挑战也是机遇。挑战在于价格竞争加剧，机遇在于中国企业在海外建厂的需求可能增加——这恰恰是越南、印尼等国吸引外资的重要来源。完整报告对不同经济体的出口结构分解，以及对“中国放缓-产能外溢”传导链条的量化分析，是理解这一复杂博弈的关键。

这是报告最具有洞察力的部分。随着能源价格下跌但美元走强，亚洲央行面临的压力正在发生轮动。报告将亚洲央行分为两类：

第一类是对美元走强最脆弱的央行。印尼央行首当其冲——印尼盾面临的压力使其在JPM已有25个基点加息预测的基础上，如果美元持续走强，加息风险将进一步上升。韩国央行同样面临风险：美元走强叠加持续的科技动能，增加了其在未来一年已经预期加息100个基点的基础上进一步加息的概率。台湾央行也可能在年底前加入加息行列。

第二类则更关注国内增长与通胀动态。对于菲律宾央行，能源价格下降意味着其终端利率可能低于JPM当前预测的6\%。对于印度央行、泰国央行和新加坡金管局，能源价格下降增强了它们按兵不动的判断—

[... middle omitted ...]

完整报告中关于各国央行反应函数的具体参数设定和情景分析，是理解这种分化的底层逻辑。

报告对中国5月财政数据给出了一个看似矛盾但逻辑自洽的解读：弱数据本身为下半年留下了更大的政策空间。5月财政支出连续第二个月收缩，基础设施支出同比下降12\%，政府性基金收支因土地销售进一步疲软而大幅下降。财政存款升至过去十年5月第二高水平，地方政府专项债券发行在6月已有所反弹。

JPM据此维持对下半年增长节奏的判断：三季度GDP季调后年化增速3.5\%，四季度3.7\%，全年增速4.7\%。这一预测建立在一个“更好但依然不均衡”的增长组合上——财政执行应会改善，全球工业产出强劲和AI资本支出周期将继续支撑中国出口和高科技/装备制造，部分抵消消费疲软。

报告还给出了一个关键的政策触发点：如果7月发布的6月活动数据非常疲弱，7月下旬政治局会议后的政策响应将更加明显——包括加快项目审批、加速地方政府专项债券资金使用、更积极动用财政存款、扩大政策性银行融资，以及对基础设施和住房稳定化的定向支持。

\subsection{[R002] GS：大宗商品真正的故事不在能源，而在金属与电力}
{\small\color{refgray}归类：能源与大宗：从能源转向金属与电力}

当霍尔木兹海峡的航道重新开放、原油价格开始回落，市场对大宗商品的注意力似乎正在消退。但GS最新的大宗商品策略报告给出了一个截然相反的判断：大宗商品配置的价值不仅没有降低，反而在结构上变得更加重要——只是主角正在从传统的油气转向电力、铜、锂和铝。

这份由Samantha Dart领衔的报告，在分析完霍尔木兹冲突的影响后，明确提出了一个核心主张：战略投资组合依然需要大宗商品，但理由正在发生变化。不再只是对冲通胀，而是应对一个更复杂的、由结构性供给约束和地缘政治重塑共同驱动的新格局。

1. 霍尔木兹冲突揭示了一个关键事实：能源市场的韧性并未改变长期脆弱性

霍尔木兹海峡中断16周，原油和成品油价格一度分别上涨43\%和63\%，欧洲天然气和亚洲LNG涨幅达到50\%-70\%。但最终，全球市场比预期更具弹性——中国在冲突首两个月大幅削减LNG进口，至今仍压低原油进口量，这帮助限制了全球市场的紧张程度。

GS团队指出，这恰恰是值得警惕的地方。市场之所以能消化这次冲击，是因为需求端主动收缩。但一旦需求反弹——他们预计冲突后原油需求将出现超过100万桶/日的战略储备补库推动——供需平衡将重新收紧。成品油市场（汽油、柴油）的回调速度将明显慢于原油。

KC评论： 报告在这里埋了一个关键伏笔：这次冲突的“软着陆”并非因为供给端弹性大，而是因为需求端主动“让路”。一旦需求恢复，能源市场的结构性紧张就会重新浮现。完整报告里有一张图表专门拆解了成品油供需平衡的预期路径，值得细看。

更重要的是，这次冲突对全球GDP的冲击被控制在了0.4个百分点，但如果冲突持续，GS估算的冲击幅度将达到2个百分点。这个数字本身就在提醒：能源市场的尾部风险并没有消失，只是暂时被掩盖了。

报告对铜的判断可能是最具前瞻性的部分。GS团队认为，铜的需求结构正在发生根本性变化——到2030年，电网和电力基础设施将贡献超过60\%的铜需求增量，叠加国防、电动车、可再生能源和AI数据中心的直接拉动。

这不是一个简单的需求增长故事。关键在于，这些需求来源具有“战略性”特征：西方老化的电网升级已被视为国家安全优先事项，因为它同时关系到AI竞争力与能源安全。这意味着铜需求对经济放缓和高价格的敏感度正在下降，与传统需求来源（如建筑和白色家电）形成鲜明对比。

供给端则面临越来越深、品位越来越低、矿石越来越硬的结构性挑战。GS预计铜需求增速将持续超过矿山供给增速，价格必须上涨才能通过刺激废铜回收和替代材料（铝）来弥合缺口。

KC评论： 报告的核心洞察不是“铜价会涨”，而是“铜价必须涨才能让市场出清”。这是两种完全不同的叙事框架。前者是预测，后者是结构性判断。完整报告里有两张关键图表——需求增速vs供给增速的对比、以及政策风险导致的铜流动态势——直接支撑这个判断。

GS因此上调了2026年底和2027年平均LME铜价预测至13,735美元/吨和13,800美元/吨，并认为到2035年需要15,000美元/吨才能维持矿山运营、提升废铜回收率并支持新矿开发。

黄金自2022年以来已上涨123\%，但GS依然看多。其2026年底目标价为4,900美元/盎司。核心驱动力是新兴市场央行的结构性多元化——自2022年俄罗斯储备被冻结以来，这一趋势就已成为黄金需求的中流砥柱。

世界黄金协会的最新调查支持这一判断：接受调查的76家央行中，创纪录的45\%预计在未来12个月内增加黄金储备，约90\%预计全球储备将继续上升。GS据此假设2026年和2027年央行将继续以每月50吨和40吨的速度增持。

但报告也坦诚地指出了短期逆风：鹰派美联储正在削弱“货币贬值”主题，市场对通胀担忧下的加息预期压制了ETF需求。GS认为这些逆风将随时间逐步消退——他们的经济学家预计2027年下半年将开启延迟但持

[... middle omitted ...]

存持续下降。周期性商品（石油和工业金属）表现最强，因为供需缺口直接推高价格。

第二种是“供给中断通胀”。霍尔木兹冲突就是典型——通胀上升、增长放缓，股债双杀。此时广谱商品（不含贵金属）是最有效的对冲，因为它们是直接被中断的投入品。

第三种是“机构信誉风险通胀”。当市场对央行或财政可持续性产生怀疑时，黄金是唯一有效的对冲。但报告明确警告：在前两种场景下，黄金不仅无效，反而可能率先下跌——因为通胀推高加息预期，提高了持有无收益资产的机会成本，同时股市下跌可能触发保证金清算，而黄金的高流动性使其成为最方便的变现工具。

这个框架的价值在于：它帮助投资者避免用黄金去对冲不该由它承担的风险。

5. 报告尚未完全回答的关键问题：电力市场的“瓶颈”何时会引爆价格

报告在结尾部分提出了一个极具前瞻性的判断：伊朗冲突最终强化的是电力和金属需求的主题，而非油气。从电动车加速替代、可再生能源投资增加、电网升级需求，到国防支出扩DaiwaAI竞赛，这些都在支撑铜、锂、铝和电力的需求。

\subsection{[R003] 摩根斯坦利：MS：融资市场正发出二季度末最危险的信号}
{\small\color{refgray}归类：风险偏好：融资市场压力与财政风险冲击}

当美联储还在通过点阵图传递“一次加息”的鹰派信号时，华尔街最敏感的角落已经亮起了红灯。MS利率策略团队在6月27日发布的一份报告中直言：股票融资市场正在经历一场堪比2024年12月年末的紧张局面，二季度末发生去杠杆事件的风险正在上升。

这份报告的核心判断并非来自传统宏观经济模型，而是来自一个多数投资者容易忽视的角落——AXW期货。这个捕捉标普500总收益期货与浮动利率之间价差的工具，正在发出近十年来最极端的信号。

对于任何一个管理着数十亿美元资产、或正在为下半年配置做决策的投资者来说，理解这个信号的含义，可能比关注美联储下一次会议的点阵图更为紧迫。

KC评论： MS不是在预测经济衰退，而是在警告一个技术性的、但可能引发连锁反应的流动性事件。历史上，这类信号往往出现在市场转折点之前。完整的报告里有一张图显示，当AXW期货超过2个标准差时，标普500在随后三个月内都会出现显著回调。这张图表值得每个人仔细看一遍。

报告揭示了一个令人不安的分化：国债回购市场目前定价相当温和，隔夜GC回购报价仅为3.75/3.72，与2025年6月季度末和今年4月纳税日时的情况类似。这意味着美联储的利率走廊并未受到实质性挑战，SOFR利率大概率会稳定在目标区间内。

然而，在股票融资的另一端，情况截然不同。7月到期的1个月AXW期货合约在上周五飙升至+200个基点，正在迅速逼近2024年12月年末前创下的历史最高水平。报告特意使用了10日移动平均线来平滑合约滚动的影响，即便如此，读数依然触目惊心。

这种分化的本质是什么？国债融资市场受益于美联储自2025年12月启动的准备金管理购买操作，流动性相对充裕。而股票融资市场则面临一个更结构性的问题：对冲基金对杠杆化多头敞口的需求，与交易商资产负债表容量之间的失衡正在加剧。

几个因素叠加推高了AXW水平：股票发行相关的资产负债表需求、杠杆ETF的增长、以及居高不下的期货持仓。但报告特别强调了一个更危险的变量：过去一周半导体板块的财报引发的市场反应，恰恰发生在对冲基金杠杆化多头敞口最集中的领域。

KC评论： 这里的关键不是某个板块的涨跌，而是“集中度”本身。当大量杠杆资金集中在少数几个板块时，任何一次波动都可能触发连锁反应。报告没有展开的是，这种集中度到底有多高、哪些基金最脆弱。这些细节正是完整版研报最有价值的部分。

MS团队用了一个精妙的框架来解释当前的风险。他们指出，当股票融资成本达到统计上的极端水平时（超过252日滚动窗口的2个标准差），标普500往往会向均值回归。这不是一个偶然的相关性，而是有因果逻辑的。

逻辑链条如下：融资成本上升 -> 杠杆投资者维持多头头寸的门槛提高 -> 部分头寸被迫平仓 -> 价格下跌 -> 更多杠杆头寸触及风险阈值 -> 进一步平仓。在这个过程中，曾经推动资产价格上涨的杠杆力量，开始向相反方向切割。

报告特别提到，这种“去杠杆事件”的风险在二季度末尤为突出。因为临近季度末，主经纪商为了避免进入下一个GSIB附加费区间，会变得更加保守——这意味着更少的融资可用性和更高的定价。对于对冲基金来说，这意味着他们可能无法在季度末之前获得足够的融资来维持现有头寸。

更关键的是远期合约的信号。报告指出，如果6月30日季度末之后，远月AXW合约（如AXW3）的融资成本没有正常化，那就意味着股票融资成本将持续高企，直到发生一次相当显著的去杠杆事件。

KC评论： MS实际上在说：市场可能正在等待一个“催化剂”。这个催化剂可能是某个大型基金的爆仓，也可能是某个宏观数据的意外。完整报告里有一张图展示了AXW3与标普500的统计关系，当AXW3超过2个标准差后，标普500在三个月内都会跌到1个标准差以下。这个时间窗口值得高度警惕。

如果去杠杆事件真的发生，它对宏

[... middle omitted ...]

紧，美联储实际上已经不需要再主动加息。报告进一步指出，如果市场开始意识到这一点，那么对美联储加息尾部风险的保险溢价就会下降，从而推动收益率曲线扁平化。

KC评论： 这是整份报告最精妙的分析。MS实际上在说：美联储可能不需要真的加息，市场已经帮它加了。如果去杠杆事件发生，金融条件会进一步收紧，届时市场会开始定价“美联储政策错误”的风险——即美联储在点阵图上暗示的加息，可能根本不会发生。报告里有一张图展示了不同美联储资产负债表制度下准备金与SOFR-IORB利差的关系，这张图对于理解当前金融条件的真实位置至关重要。

尽管MS的分析框架非常清晰，但报告本身也留下了一些未完全展开的问题。

第一个问题是：去杠杆的具体触发点在哪里？报告提到了半导体板块的财报，也提到了对冲基金在AI相关领域的集中度，但没有明确量化“什么水平的融资成本会触发系统性平仓”。AXW期货目前是+200bp，2024年12月的高点是多少？如果继续上行到+250bp或+300bp，哪些类型的基金最脆弱？

\subsection{[R004] 美国银行：存储器超级周期可能延续到2030年}
{\small\color{refgray}归类：AI与算力：存储器超级周期与AI落地组织变革}

当多数人还在讨论半导体周期何时见顶时，美国银行最新研报给出了一个反直觉的判断：存储器行业的超级周期不仅没有结束，反而可能延续到2027年甚至2030年。这份基于美光科技最新财报和指引的分析，揭示了一个正在发生的结构性转变——存储器正从周期性行业转向类似台湾晶圆代工的高利润、低波动模式。

报告的核心主张是：存储器行业正在经历一场由长期协议(Long-Term Agreements, LTA)驱动的结构性变革，这将从根本上改变行业的周期特征和盈利模式。这不是又一个周期高点，而是一个新范式的起点。

美光科技最新季度财报给出了几个令人瞩目的数字：5月季度营收达到410亿美元，同比增长346\%；毛利率85\%，营业利润率81\%；对8月季度营收指引为500亿美元，意味着环比增长21\%。这些数字本身已经足够惊人，但更关键的是美光管理层对行业前景的展望。

美国银行分析师从中提炼出五个对亚洲存储器厂商具有深远影响的判断：

第一，超级周期或芯片短缺可能持续到2027年甚至2030年。这不是一个短期现象，而是由AI需求驱动的长期结构性变化。

第二，与大型科技公司和OEM的长期协议(LTA)正在增加，这有望使行业变得不那么周期性，并实现高利润——类似于台湾晶圆代工的模式。

第三，新建晶圆厂并非易事，面临高建设成本、地方政府监管、电力/水供应问题等多重挑战。

第四，即使到2028年，晶圆产能扩张也有限，因为需要大量洁净室用于旧厂升级，制造周期更长，前端甚至后端设备尺寸更大，先进芯片良率较低，以及HBM与常规DRAM之间较高的转换比例。

第五，尽管资本支出相比正常周期翻倍以上，自由现金流仍在显著增加。

KC评论： 这五个信号中最关键的是LTA的普及。如果存储器行业真的向晶圆代工模式转型，那么投资者需要重新评估整个行业的估值框架。周期性行业的低估值逻辑将不再适用。完整报告中包含了对这些信号更详细的拆解，以及每个信号对不同厂商影响的具体分析。

2. LTA正在重塑存储器行业的盈利结构，但并非所有厂商都能受益

美国银行指出，美光和其他大型存储器厂商越来越依赖LTA来固定ASP，同时其芯片生产将更加高端化（HBM、SOCAMM、LPDDR5）。这为存储器厂商开辟了新的增长机会，而无需面临激烈的价格竞争。

这里的关键在于，LTA改变了存储器行业的定价机制。传统上，存储器价格由现货市场供需决定，波动剧烈。但LTA意味着价格被提前锁定，供需双方共同承担风险，这降低了价格波动的幅度和频率。

然而，并非所有厂商都能享受这种新范式的好处。报告特别指出，南亚科技（Nanya Tech）的DRAM销售（主要是旧款/传统DRAM）仍然主要基于月度/季度协商的ASP，对LTA的敞口非常低。

KC评论： 这意味着存储器行业内部正在发生分化。能够获得LTA的厂商将享受更稳定的收入和利润，而那些依赖现货市场定价的厂商将面临更大的不确定性。完整报告中对不同厂商的LTA敞口和产品结构有更详细的对比分析，这是判断投资标的的重要维度。

3. 2026年DRAM营收将接近翻四倍，但增长节奏正在变化

美国银行更新了全球存储器预测，预计2026年全球DRAM营收将同比增长316\%，NAND营收同比增长295\%。这些数字背后是ASP的强劲反弹：DRAM ASP预计同比增长242\%，NAND ASP同比增长234\%。

但报告同时指出一个重要的变化：ASP可能不再强劲上涨，如果基于LTA的销售增加。美光8月季度营收环比增长约20\%，远低于2月/5月季度的75\%/74\%环比增速。这意味着增长正在从“价格驱动”转向“量价平衡”。

美国银行小幅上调了 年全球DRAM/NAND销售预测2-4\%，但调整主要基于对2Q26 ASP的乐观预期，部分抵消了LTA导致的3Q或

[... middle omitted ...]

8\%，已经连续五个月实现三位数增长。5月出口额达到255亿美元，环比增长16\%，是2025年平均水平的三倍以上。

韩国半导体出口主要由存储器构成，这一数据与美光的乐观指引高度一致。同时，DRAM现货价格已连续五周上涨，累计涨幅20\%，而NAND现货价格则下跌5\%。

KC评论： DRAM和NAND现货价格的分化值得关注。DRAM的强势反映了HBM和高端DRAM的需求旺盛，而NAND的疲软可能意味着供应端压力或需求结构不同。完整报告中对不同存储器产品的供需分析更加细致，包括HBM4、SOCAMM、LPDDR5、GDDR7、eSSD等具体产品的需求和产能情况。

尽管美国银行对存储器超级周期给出了乐观判断，但报告中仍有一些未完全回答的问题。

第一个问题是产能扩张的实际边界。报告指出“即使到2028年晶圆产能扩张也有限”，但给出了DRAM产能2026年同比增长11\%、2027年增长9\%的预测。这个增速在历史上并不算低。真正的约束可能不是产能总量，而是先进制程产能的比例。

\subsection{[R005] Bernstein：钠离子电池的“奇点”不在车，在储能}
{\small\color{refgray}归类：能源与大宗：从能源转向金属与电力}

全球储能行业正在经历一个被多数人低估的技术拐点。过去几年，市场焦点始终落在锂价的剧烈波动和电动车渗透率的爬坡节奏上，但一份来自Bernstein的最新周报揭示了一个更值得关注的信号：钠离子电池正在从实验室验证走向商业化部署，而它的第一个引爆场景不是电动车，而是大型储能系统。

这份报告汇集了过去一周全球储能产业链的数十条动态，从CATL推出全球首个公用事业级钠离子储能系统Tener Sodium，到中国钠创科技在宁夏启动10GWh钠离子电池项目，再到Hina电池与一汽解放在极寒条件下完成钠离子重卡电池测试——这些事件拼在一起，指向同一个判断：钠离子电池在储能领域的成本拐点可能比市场预期的来得更早。

对于关注能源转型、电力系统投资和电池供应链格局的读者来说，这是一个需要重新校准认知框架的时刻。

KC评论： 市场对钠离子电池的普遍认知还停留在“能量密度低、只能用在低端两轮车”的阶段。但Bernstein这份周报的拼图显示，钠离子电池的真正突破场景是储能——在那里，能量密度不是核心矛盾，成本、安全性和循环寿命才是。这个认知差本身就是投资机会。

1. CATL的Tener Sodium不是概念产品，而是商业交付的起点

CATL在本周正式发布的Tener Sodium，是全球首个面向公用事业级的钠离子储能系统。这不是一款还在PPT阶段的概念产品。报告明确提到，该系统支持单机超过30MWh的容量，储能时长可在1至8小时之间灵活配置，循环寿命达到15,000次，且在极端温度下表现稳定。更重要的是，CATL已经确认技术和供应链均达到商业化水平，计划2026年9月开始在中国交付，2027年推向全球市场。

这些参数意味着什么？储能时长1到8小时的灵活性，直接对应了电网侧调峰、调频以及工商业削峰填谷的核心需求。15,000次的循环寿命，意味着在每日一充一放的场景下，系统可以使用超过40年——远超当前锂电储能系统的设计寿命。而极端温度下的稳定性，则解决了锂电在高温和低温环境中性能衰减的痛点。

Bernstein没有在报告中直接给出Tener Sodium的度电成本，但从这些技术指标可以推断：如果CATL能够在2026年如期交付，这将意味着钠离子储能的度电成本已经具备与磷酸铁锂储能竞争的能力。

KC评论： 这里有一个关键问题报告没有完全回答：Tener Sodium的实际成本与磷酸铁锂相比到底如何？CATL没有公布具体价格，但15,000次的循环寿命和极宽的工作温度范围，意味着系统在全生命周期内的度电成本可能已经低于当前主流的LFP储能方案。这个假设需要完整报告中的详细技术参数和成本模型来验证。

CATL的Tener Sodium并非孤立事件。Bernstein周报中同时记载了另外两条重要动态：中国钠创科技在宁夏启动了10GWh钠离子电池项目的建设，总投资30.3亿元人民币，覆盖从正负极材料到电芯制造、PACK组装和储能应用的全链条。与此同时，Hina电池与一汽解放在零下40摄氏度的极寒条件下完成了339kWh钠离子重卡电池的7个月、15,000公里路测，结果显示电池可用容量保持率超过90\%，支持20至25分钟快充，循环寿命超过8,000次。

这三条信息叠加在一起，描绘出的图景是：中国正在从材料端、制造端到应用端，全面推动钠离子电池的产业化。宁夏项目代表了供应链的规模化落地，Hina的路测证明了钠离子在商用车领域的可行性，而CATL的Tener Sodium则标志着钠离子正式进入公用事业级储能市场。

对于全球储能产业链来说，这意味着一个重要的结构性变化：当中国以每年数十GWh的规模推进钠离子电池产能建设时，全球钠离子电池的供应链成本曲线将沿着与光伏、锂电相似的路径快速下降。任何依赖单一技术路线的储能投资策略，都需要重

[... middle omitted ...]

Northvolt在德国海德的未完工工厂，但产能规模远小于Northvolt原计划。

这些案例共同指向一个结论：欧洲电池供应链的自主化进程，远没有市场预期的那么顺利。监管审批迟缓、融资环境收紧、技术人才短缺，叠加中国企业的成本优势和规模效应，使得欧洲本土电池制造的竞争力在短期内难以形成。

对于投资者而言，这意味着两个推论：第一，欧洲储能市场的电池供应在中期内仍将高度依赖中国进口，关税和碳边境调节机制虽然会增加成本，但不会改变这一基本格局；第二，那些在欧洲拥有本地化产能的中国电池企业，如CATL和比亚迪，将在欧洲储能市场中占据更有利的竞争位置。

KC评论： 报告没有展开的一个关键问题是：欧洲电池本土化的失败，是否意味着中国电池企业将面临更大的地缘政治风险？从目前的信息看，欧洲的态度是矛盾的——一方面希望减少对中国供应链的依赖，另一方面又离不开中国企业的技术和产能。这个矛盾的演变方向，将直接影响中国电池企业的估值逻辑。完整报告中的竞争格局分析或许能提供更清晰的框架。

\subsection{[R006] GS：韩国军工和机器人，下半年最值得关注的两个板块}
{\small\color{refgray}归类：区域市场：韩国军工与机器人、人民币汇率信号}

当全球投资者的注意力仍被AI和半导体垄断时，一份来自GS的最新调研笔记揭示了一个被忽视的信号：韩国军工和机器人板块，正在酝酿一轮由订单、政策和估值共同驱动的修复行情。

这份报告来自GS韩国工业团队，基于上周在纽约和波士顿的路演反馈写成。核心判断清晰而直接：韩国军工和机器人板块，有望在下半年迎来显著表现。

这不是一个关于“长期趋势”的模糊论断，而是一个关于“催化剂正在密集兑现”的即时判断。报告给出了四个关键支撑：估值压缩后的安全边际、订单加速的可见性、欧洲市场敞口的扩大，以及全球政策层面对机器人产业的集体押注。

对于正在寻找低相关性、有事件驱动逻辑的产业投资者而言，这份报告提供了一个值得认真对待的观察框架。

GS指出，韩国军工股在二季度显著跑输大盘，原因有三：订单进展停滞、市场资金向存储芯片集中、以及全球军工同行股价走弱。

过去一个月，韩国军工行业出现了至少五个关键催化事件：瑞士因爱国者导弹交付延迟而转向韩国防空系统；西班牙即将敲定韩华航空航天K9自行榴弹炮订单；LIG D\&A与莱茵金属签署防空导弹谅解备忘录；泰雷兹与韩华航空航天就Chunmoo制导导弹达成合作；秘鲁总统选举尘埃落定，此前因地缘政治不确定性而搁置的Rotem订单障碍正在解除。

KC评论： 这些事件的意义不在于单个订单的金额，而在于它们共同指向一个结构性变化——西欧国家正在将韩国视为传统军工供应链之外的“第二选择”。当交付速度成为关键变量时，韩国的产能和执行力就构成了核心议价权。报告中的估值图表显示，韩华航空航天和现代Rotem的12个月前瞻市盈率已处于近年低位，而盈利预期仍在上升。这种“估值压缩+盈利增长”的组合，往往是机构资金重新配置的起点。

如果说军工板块的逻辑是“订单驱动估值修复”，那么机器人板块的逻辑则是“政策驱动预期重塑”。

GS在报告中强调，机器人已经成为全球性的政策优先事项。这不是一个空洞的判断。报告列举了三个具体信号：韩国科技部正在建立机器人数据训练中心，政府资金已列入明年预算；韩国总统将于6月29日召集主要机器人CEO举行圆桌会议；美国商务部长Lutnick将机器人定义为“技术竞争的下一个战场”，并将中国的国有支持产业视为国家安全威胁。

这些信号叠加在一起，意味着机器人产业正在从“创业公司讲故事”的阶段，进入“国家战略资源配置”的阶段。

KC评论： 投资者常问的一个问题是：现在布局机器人是否太早？GS的回应值得细读：这不是一个“时间点”的问题，而是一个“竞赛”的问题。当主要经济体都在加速投入时，等待“完美时机”可能意味着错过定价窗口。报告中引用的风险投资数据也佐证了这一点——机器人领域的VC交易活动已达到历史最高水平。

GS特别提到了英伟达推出的Halos——一个开放的机器人安全系统。这个看似技术细节的事件，实际上具有产业里程碑意义。

当前人形机器人商业化的最大障碍之一，是安全认证的碎片化。每家厂商需要独立完成漫长的认证流程，这推高了成本、拉长了周期。Halos的作用在于标准化，它将帮助行业参与者加速安全认证，从而降低商业化的准入门槛。

报告还引用了一个更宏观的观察：基于强化学习的AI应用，在LLM领域已经展现出持续的能力突破曲线。机器人的能力提升，很可能遵循类似的轨迹。

KC评论： 报告中的图表显示，AI能够完成的任务时长（以50\%成功率为标准）正在以指数级增长。如果机器人领域复制这一路径，那么当前看似遥远的“人形机器人普及”可能比市场预期来得更快。GS推荐的具体标的覆盖了现代汽车集团相关公司（如现代摩比斯）、汽车零部件公司（如HL Mando）和纯机器人公司（如Robotis），这种“从零部件到整机”的覆盖逻辑，反映的是对整个产业链的押注。

尽管GS的判断清晰且具有说服力，但这份报告也留下了一些

[... middle omitted ...]

目标PE为60倍，基于2035年盈利折现。这种长期折现模型本身隐含了对“人形机器人大规模商用”的极高置信度。如果商业化进程慢于预期，或者关键订单（如人形机器人执行器）在 年未能兑现，估值面临的风险不容忽视。

这些未解问题，恰恰是深度阅读完整报告的价值所在。GS的原始报告包含了更详细的估值图表、盈利预测拆分，以及每只推荐股票的具体风险分析——这些细节在摘要中无法充分展开。

每天，我们的AI agent会结合人工审核，生成国际顶级投行的中文摘要与深度评论，覆盖GS、MS、JPM、UBS、Citi等机构的最新研报。更新频率为每日10-40页，包含当天最新数据图表合集。这些内容既方便直接喂给AI模型进行二次分析，也方便人工快速把握全球市场的核心动态。欢迎加入我们的星球微信群，继续讨论这些未解问题，获取完整报告的原始图表与详细拆解。

Personal reading notes and learning share only. Not investment advice.

\subsection{[R007] NOM：人形机器人量产拐点已至，但真正的主线不在机器人本身}
{\small\color{refgray}归类：人形机器人：量产拐点与反射层架构}

人形机器人行业正在经历一个被市场低估的转折。NOM在最新发布的研报中给出了一组值得深思的数字：特斯拉Optimus的产能规划从2027年的5万台上调至7万台，2028年再增加一条产线，远期总产能向150万台迈进。但真正重要的不是这个数字本身，而是支撑这一数字背后正在发生的两个结构性变化——数据采集方式的根本性转向，以及整机代工模式的加速成型。

这份报告基于2026年6月中旬对宇树科技、Deep Robotics、申昊科技、梅卡曼德以及特斯拉供应链企业的实地调研，其核心判断是：人形机器人的量产逻辑正在被重写。过去行业争论的焦点是“谁能造出更好的硬件”，而NOM认为，未来12个月的决定性变量是“谁能以更低成本获取更高质量的训练数据”以及“谁能在代工链条中占据不可替代的位置”。

这不是一个关于技术突破的故事，这是一个关于成本曲线和供应链权力转移的故事。

1. 2026年出货量预期大幅上修，消费端而非工业端成为主要驱动力

NOM预计2026年中国市场的人形机器人出货量将达到4万至5万台，较此前预期显著上调。这个数字背后有两个关键驱动因素：一是政府主导的“具身智能基地”采购加速，二是预计2026年下半年由低价产品发布引发的消费需求拐点。

从下游结构来看，NOM的估算显示，消费端和表演/娱乐端各占约30\%，政府采购（数据采集）占20\%，教育占15\%，商业和工业应用仅占3\%至5\%。这意味着，当前人形机器人的主要出货场景仍然是“数据采集”和“体验消费”，而非真正意义上的工业替代。

KC评论： 工业场景占比仅3\%至5\%，这个数字值得反复咀嚼。市场对机器人的想象往往集中在工厂流水线，但现实是，工业部署在精度、节拍时间和硬件成本上仍未通过验证。报告明确指出，工业落地仍处于“验证瓶颈”中。这意味着，短期内的投资逻辑不应建立在“机器人替代工人”的宏大叙事上，而应关注消费端和政府端的需求能否持续。

NOM进一步给出了中国厂商的分层出货预估：两家头部企业（宇树科技和另一家未具名厂商）各自出货1万至1.5万台，同比增长2至3倍；第二梯队约3000台；第三梯队500至1000台。这个金字塔结构本身就揭示了行业的竞争格局——头部集中度极高，且正在加速拉开差距。

2. 数据采集范式正在发生根本性转变：非机器人数据成本仅为机器人数据的20\%

这是整份报告中最具洞察力的判断之一。NOM通过产业调研发现，行业正在加速从“真实机器人数据采集”转向“非机器人数据采集”，包括远程操作、通用操控接口和第一人称视角数据。后者的成本仅为前者的约20\%，且采集速度显著更快。

NOM预计，混合数据训练（约90\%非机器人数据+10\%真实机器人数据）将成为2027年的标准做法。这意味着，2026年可能是真实机器人数据采集需求的峰值年份，随后将逐年下降。政府目前仍是数据采集工厂的主要出资方，一个约1000台规模的大型数据工厂需要投资5000万至1亿元人民币，投资回收期为3至5年。

KC评论： 这个转变的意义被市场严重低估。如果非机器人数据能够有效替代真实机器人数据，那么人形机器人OEM的硬件采购压力将大幅下降——它们不再需要为了采集数据而大规模生产机器人。但这也提出了新的问题：数据质量、一致性和跨平台标准化如何保障？NOM在报告中承认“这些门槛正在提高”，但并未给出量化的验证。这正是完整报告中最值得细读的部分——它包含了对数据质量评估框架的初步探讨。

NOM的调研显示，硬件供应链仍然是资本密集型且投资回收期长的领域，这迫使大多数二线和三线OEM选择外包整机生产。长期来看，外包比例可能进一步上升。但也有一些例外，比如宇树科技仍然坚持自建供应链，管理层认为这可以降低整机物料清单成本，直接塑造定价能力。

在代工企业中，NOM判断，汽车供应链企业可能拥有更强

[... middle omitted ...]

力；另一方面，这也意味着竞争正在快速从“技术竞赛”转向“成本竞赛”。如果代工模式成为主流，那么真正的价值创造者可能不再是机器人品牌商，而是那些能够同时掌控供应链管理、大规模制造和成本控制能力的代工厂商。NOM提到汽车供应链企业可能更具优势，但并未展开论证——这些企业的利润率结构、客户粘性和技术壁垒如何？完整报告中有更详细的竞争格局分析。

4. 报告尚未完全回答的关键问题：软件迭代能否突破工业验证瓶颈

NOM在报告中多次提到，更快的软件迭代正在解锁零售、导览和展览等场景的应用，但工业部署仍然受限于精度、节拍时间和模型训练成本。这是一个“尚未完全回答”的关键问题。

具体来说，NOM认为工业场景的验证瓶颈包括：精度和节拍时间尚未满足要求，模型训练成本仍然过高。但报告没有给出这些瓶颈何时可能被突破的判断，也没有分析不同工业场景（如汽车装配、物流分拣、精密制造）的差异化验证门槛。这是读者在阅读完整报告时需要重点关注的部分——NOM在附录中提供了更详细的场景拆解和成本模型。

\subsection{[R008] NOM：人民币中间价模型透露的政策信号，比点位更重要}
{\small\color{refgray}归类：区域市场：韩国军工与机器人、人民币汇率信号}

2026年6月29日，NOM亚洲外汇策略团队发布了一份看似技术性的报告——USD/CNY固定汇率模型更新。报告本身聚焦于模型投影值的微调，但在当前中美关系、国内经济政策博弈的关键节点，这份报告透露出的信息远比一个数字变化更有深意。

*核心判断：人民币中间价的形成机制正在经历一个从“被动跟随”到“主动引导”的微妙切换。** 模型投影值与实际中间价的偏离幅度、逆周期因子的调整方向，正在成为观察中国央行汇率管理意图的前瞻指标。对于高净值投资者和产业决策者而言，理解这些技术细节背后的政策逻辑，比猜测具体点位更重要。

报告给出的最新模型投影值为6.8077，较前次的6.8166下降了89个基点。计入逆周期因子后的投影值为6.8120，较前次下降46个基点。这两个数字本身并不惊人，但变化的方向和幅度，结合近期政策事件的时间表，构成了一个值得深度拆解的信号矩阵。

NOM模型的核心输出是USD/CNY中间价的投影值。6月29日的数据显示，投影值从6.8166降至6.8077，降幅为89个基点。这意味着，基于模型捕捉的隔夜市场变量——包括美元指数、欧元/美元、美元/日元等主要货币对的变动——模型认为央行应将中间价设定在更低的位置。

这一变化本身并不意外。过去几周，美元指数因美国经济数据走弱而出现回调，人民币的贬值压力也随之有所缓解。但关键在于，模型投影值的下降幅度大于市场预期，暗示外部压力正在边际减弱。

KC评论： 对于普通投资者而言，不必纠结于模型的具体公式。关键信号是：当模型投影值持续下降，意味着央行面临的“被动贬值”压力在减轻，这为其主动管理汇率提供了更大的腾挪空间。完整报告中包含了模型误差的历史图表，这些数据能帮助判断模型预测的可靠性——值得仔细看。

2. 逆周期因子调整幅度收窄，央行正在减少对市场的“暴力干预”

模型投影值与计入逆周期因子后的投影值之间的差距，是观察央行干预意愿的直接窗口。本次报告显示，计入逆周期因子后的投影值为6.8120，与无因子投影值6.8077的差距仅为43个基点。而前次报告的差距为50个基点（6.8166 vs 6.8120）。

差距的收窄意味着：央行正在减少逆周期因子的使用力度。这有两种可能的解读：一是央行认为当前市场运行相对平稳，无需过度干预；二是央行希望在关键政策事件（如7月政治局会议、APEC峰会）前，保留更多的政策工具空间。

KC评论： 逆周期因子是央行调控中间价的“隐形之手”。当这个因子使用力度加大时，往往意味着央行在对抗市场单边预期；当力度减弱时，则可能表明央行对市场自我修复有信心。报告中的模型误差图表（Fig. 2）展示了历史误差分布，可以帮助判断当前干预力度处于什么历史分位——这是完整报告中最值得深挖的数据之一。

3. 7月政治局会议是近期最重要的政策观察窗口，汇率管理或成议题

报告在日历部分列出了2026年下半年的一系列关键事件，其中7月下旬的政治局会议被放在首位。按照惯例，7月会议通常会对下半年经济工作做出部署，其中汇率政策、资本流动管理、货币政策取向等，都可能成为讨论重点。

当前中国面临的内外环境复杂：国内经济复苏基础尚不牢固，房地产市场仍在调整；外部中美关系虽有缓和迹象，但关税、科技封锁等结构性问题未解。在此背景下，央行如何在“保持汇率弹性”和“防范汇率超调”之间取得平衡，将成为市场关注的焦点。

报告虽然没有直接给出政策建议，但通过模型投影值与实际中间价的对比，实际上为读者提供了一个“预判央行意图”的框架：如果模型投影值与实际中间价的偏差持续扩大，意味着央行在主动引导汇率；如果偏差收窄，则意味着央行更倾向于让市场定价。

报告日历中列出了两个重要事件：11月的APEC峰会（中国将在深圳主办）和年底中国国家主席对美国的访问。这两个事件的时间节点

[... middle omitted ...]

风险管理；对于投资者而言，这意味着需要将政治事件纳入汇率定价模型。

5. 报告尚未完全回答的关键问题：模型与实际中间价的偏离，是主动引导还是被动误差？

NOM模型虽然提供了清晰的投影值，但报告并没有详细解释模型误差的来源。从Fig. 2的模型误差图表来看，历史误差有时相当大——这意味着模型并非完美预测工具。当实际中间价与模型投影值出现较大偏差时，究竟是央行在主动引导，还是模型本身存在局限性？

这个问题没有标准答案，但值得深入探讨。如果偏差是主动引导的结果，那么投资者应该关注央行引导的方向和力度；如果偏差是模型误差，那么投资者应该关注模型的修正机制和迭代逻辑。

KC评论： 这是整份报告最值得追问的地方。报告给出了模型投影值，但没有给出“实际中间价”与投影值的对比数据。要回答这个问题，需要读者自行查阅央行每日公布的中间价，并与模型投影值做对比。完整报告中包含了历史误差图表，可以帮助判断当前偏离是否在正常范围内——这是加入社群继续讨论时，我们可以一起拆解的核心问题。

\subsection{[R009] Bernstein：斗山能源的2500亿美元核能订单，不只是“韩国制造”的故事}
{\small\color{refgray}归类：能源与大宗：从能源转向金属与电力}

Bernstein：斗山能源的2500亿美元核能订单，不只是“韩国制造”的故事

全球核电建设正在进入新一轮加速周期，但真正值得关注的不是总装机量，而是供应链中少数几家具备重型设备制造能力的公司。Bernstein最新发布的亚太核电深度报告，将目光锁定在一家韩国公司——斗山能源（Doosan Enerbility）。报告的核心判断是：在未来十年，斗山能源面临约2500亿美元的可参与核电项目总市场规模（TAM），其中约440亿美元（67万亿韩元）是可直接转化为设备合同的订单机会。这并非一个模糊的行业乐观叙事，而是一个基于项目进展节点、可追踪、可验证的订单管道。

这份报告的独特价值在于，它没有停留在“核电要复兴”的宏大判断上，而是逐国别、逐项目地拆解了哪些项目已经进入设备采购前阶段，哪些还在可行性研究，以及这些时间节点对斗山这样的设备商意味着什么。对于关注产业投资、能源供应链和制造业升级的读者，这是一份少见的、从“订单可见性”出发的深度分析。

1. 2500亿美元的可参与市场，核心不是规模而是“可执行性”

Bernstein的分析框架有一个关键前提：他们只统计那些已经推进到设备采购前阶段的项目。这意味着这些项目已经完成了或正在完成监管许可、前端工程设计（FEED）以及融资结构确定。在这个阶段之前，项目仍存在高度的不确定性，无法作为订单预测的依据。

基于这一筛选标准，报告识别出约35吉瓦（GW）的核电项目，总项目成本约2500亿美元。这并非所有在建或规划中的核电项目，而是斗山能源“真正有机会拿到设备合同”的那一部分。从区域分布看，欧洲项目是近期订单的主要来源，美国项目则要到2020年代末才会随着融资和供应链支持改善而逐步落地。

这一判断对投资者的含义是明确的：斗山能源的订单增长并非依赖宏观叙事的推动，而是有具体项目节点作为支撑。波兰、保加利亚、捷克等国的项目已经进入或即将进入设备招标阶段，这些项目的推进节奏决定了斗山未来几年的订单可见度。

KC评论： 很多研究报告在讨论核电时喜欢用“万亿级市场”这样的词，但Bernstein的做法更务实。他们问的是：哪些项目已经过了“纸上谈兵”的阶段，真正到了要下单买反应堆压力容器和蒸汽发生器的时候？这才是设备商能转化为收入的市场。完整报告里有一张非常详细的国别项目表格，列出了每个项目的建设时间表、开发商、总成本和斗山可能拿到的合同金额，值得仔细看。

核电设备制造不是一个可以靠资本快速复制的行业。斗山能源的竞争力来自其位于韩国昌原的垂直整合制造基地。这个基地几乎涵盖了从原材料炼钢到最终组装的全流程，能够生产反应堆压力容器、蒸汽发生器、稳压器、汽轮机系统等核心部件。

报告特别强调了一个数据：斗山在过去40年里，已经制造并交付了34个反应堆压力容器和124个蒸汽发生器。在中国和俄罗斯之外的全球核电设备供应商中，斗山是经验最丰富的企业。这意味着任何新的竞争者想要达到同等的制造能力和质量信誉，都需要数十年时间的持续投入和项目验证。

从技术路线看，斗山在AP1000和APR1400这两种主流大型反应堆设计上都有丰富的交付记录。这种双线布局使得它能够灵活响应不同国家、不同技术路线的项目需求，而不局限于某一种堆型。此外，斗山从五年前就开始布局小型模块化反应堆（SMR）的制造能力，虽然SMR在短期内的订单贡献有限，但为中长期增长提供了额外的期权价值。

KC评论： 核电设备制造有点像高端芯片制造——不是有钱就能快速建起来的。斗山昌原工厂的垂直整合能力，意味着它可以控制关键部件的质量和交付周期，这在核电项目动辄延期超支的背景下，是一个被低估的竞争优势。报告里有一张图对比了全球主要核电项目的建设成本和工期，斗山参与的项目在效率和成本控制上明显优于行业平均，这份“经验曲线”优势是定价权的

[... middle omitted ...]

特别指出，斗山正在扩大燃气轮机产能，预计到2030年，燃气轮机订单将从2025年的4.7万亿韩元增长到约7万亿韩元。

但报告也坦承，执行风险和项目延期是主要的不确定性。为此，Bernstein在订单预测中引入了风险调整因子，根据项目通过关键里程碑的概率对订单金额进行折让。这意味着21万亿韩元的预测并非乐观情景，而是基于概率加权后的“基准情景”。

斗山能源的利润率改善并非来自简单的规模扩张，而是业务结构的质变。报告预测，到2028年，核电和燃气轮机设备将占公司独立营收的75\%以上，而目前这一比例显著更低。这种结构变化将推动营业利润率从目前的低个位数提升到2028年的低双位数。

为什么结构变化如此重要？因为不同业务的利润率差异巨大。全球燃气轮机OEM厂商正在向20\%以上的利润率迈进，核电设备供应商的EBITDA利润率也维持在15\%左右的中高水平。相比之下，斗山传统上依赖的煤炭和海水淡化业务利润率较低且波动性大。随着高利润业务的占比提升，公司的整体盈利质量将显著改善。

\subsection{[R010] NOM：人形机器人的“脊髓”比“大脑”更关键}
{\small\color{refgray}归类：人形机器人：量产拐点与反射层架构}

当市场还在争论大模型能否赋予机器人通用智能时，一份来自NOM的最新研报提出了一个反直觉的判断：物理人工智能落地的真正瓶颈，不是语言模型的推理规模，而是能否构建一个低延迟的“反射层”——相当于人类的脊髓反应。

这份报告以NXP提出的NeuralAxis架构为理论蓝图，以宇树科技新发布的WVLA2.0模型为产品验证，勾勒出一条与当前主流“端到端大模型”路径截然不同的商业化路线。其核心主张是：在物理世界里，40毫秒的反射响应比300毫秒的深度思考更决定生死。这一判断不仅挑战了人形机器人行业的技术共识，更对工业自动化、服务机器人乃至自动驾驶的落地节奏提供了新的观察框架。

NOM分析师Frank Fan和Donnie Teng在实地调研宇树科技后，将这份研报定位为“从架构蓝图到可交付产品的完整推演”。它没有停留在概念层面，而是给出了具体的技术指标、商业化场景排序，以及当前仍未解决的量化验证问题。

Moravec悖论指出，对人类来说困难的高阶推理对机器而言相对容易，而人类无意识完成的感知与运动控制对机器却异常艰难。NOM报告借NXP的NeuralAxis框架，将这一悖论具象化为一个工程问题：物理AI系统必须在毫秒级完成对环境的反应，而不是等待云端大模型返回推理结果。

NeuralAxis将系统架构拆解为三个解耦但协同的层级：推理层（大脑皮层，约300毫秒响应）、协调层（小脑，负责运动控制与平衡）、反射层（脊髓，最低可达40毫秒）。这一架构最激进的设计在于，反射处理器被分布部署到关节、手部和足部，而非集中在一个中央大脑。这意味着，当机器人抓取一个未知重量的物体时，手指关节的本地处理器可以在40毫秒内自主调整抓握力度，无需等待上层推理指令。

KC评论： 这相当于把机器人的“本能反应”从云端或中央处理器剥离出来，下沉到执行端。对于工业场景而言，这意味着机器人可以在不增加算力成本的前提下，大幅提升对突发干扰的响应速度。完整报告中有NeuralAxis与人类神经系统对比的架构图，值得仔细研究其三层之间的通信协议设计。

这一架构的商业化含义是明确的：在制造业、物流等对实时性要求极高的场景，纯粹依赖大模型推理的机器人可能面临“反应太慢”的致命缺陷。NOM的行业调研显示，基于NeuralAxis架构的工业机器人已经在部分产线上实现了显著的生产率提升，诊断机器人的销售也在快速爬坡。

2. 宇树WVLA2.0不是VLA的简单升级，而是世界模型与动作生成的融合

如果说NeuralAxis定义了“应该做什么”，那么宇树的WVLA2.0则展示了“如何做到”。这个经过两年开发的模型，是NOM认为“第一个可商业部署的迭代版本”。它的核心差异在于，没有像多数同行那样押注单一的VLA（视觉-语言-动作）架构，而是将WMA（世界模型动作）的预测能力与VLA的端到端动作生成进行了融合。

这种融合在技术层面解决了几个关键问题：高层次任务理解、2D/3D空间语义推理、受动力学约束的动作生成，以及抗干扰能力。NOM报告详细描述了其感知系统——四路并行视觉流（RealSense深度相机、Livox MID360激光雷达、两个侧向摄像头）拼接出360度环境表征，在干扰下位置更新速度仍能控制在10毫秒以内。

更值得注意的是其“轻量化”设计。边缘计算功耗被控制在100 TOPS以内，完全运行在宇树G1 EDU搭载的NVIDIA Jetson Orin NX上，不依赖云端。这意味着即使在网络断开的情况下，机器人仍能保持完整功能。在NOM观察的演示中，一台搭载WVLA2.0的G1机器人在被干扰的会议室中，自主完成了六个连续任务，推理闭环速度约为每次90毫秒，相当于每秒执行约十次感知-预测-决策-动作循环。

KC评论： 90毫秒的闭环速度是衡量“实时性”的关

[... middle omitted ...]

的是不再依赖机器人本体在真实环境中反复试错来收集数据，而是通过仿真、合成数据或其他非物理手段生成训练素材。这一转变对于降低数据采集成本、加速模型迭代至关重要。

然而，宇树的真正护城河并不在于技术架构本身，而在于其“全栈自研+全球机队真实数据”的组合。NOM指出，宇树拥有从电机、关节到控制器的全栈自研能力，同时其全球部署的机器人机队每天都在产生海量的真实环境数据。这种“真实世界数据飞轮”是纯云端模型厂商无法复制的资产。

这引出了一个更深层的问题：在物理AI领域，算法的通用性固然重要，但谁能更快地积累真实场景下的长尾数据，谁就更有可能在商业化落地中占据先机。宇树的策略是先用工业制造场景（自有工厂的关节电机装配、上下料、夹具搬运）作为落地起点，再逐步拓展到物流分拣、柔性3C组装，最后才是家庭和医疗场景。这个顺序安排体现了对“场景难度”的清醒认知——开放、非结构化环境的复杂度呈指数级上升。

尽管NOM这份研报提供了丰富的技术细节和商业化判断，但仍有几个关键问题悬而未决。

\subsection{[R011] 国际清算银行：财政风险冲击正在制造“滞胀陷阱”，而央行宽松只会让情况更糟}
{\small\color{refgray}归类：宏观与利率：亚洲央行轮动与全球财政风险 / 风险偏好：融资市场压力与财政风险冲击}

国际清算银行：财政风险冲击正在制造“滞胀陷阱”，而央行宽松只会让情况更糟

当市场开始重新评估政府债务的安全性，后果并不仅仅是主权债券收益率上行。国际清算银行（BIS）最新工作论文揭示了一个更令人不安的传导链条：财政风险冲击会同时推高通胀和压低产出，形成典型的“滞胀”格局。更关键的是，如果央行选择用宽松货币政策来对冲，结果不是缓解，而是让通胀更高、衰退更深。

这份由BIS经济学家Denis Gorea、Ding Xuan Ng和Fabrizio Zampolli完成的第1364号工作论文，构建了一套全新的识别框架，从债券市场数据中分离出纯粹的“财政风险冲击”，并追踪其对12个经济体宏观金融变量的影响。结论清晰且具有政策警示意义：财政纪律的丧失，其代价远不止于更高的融资成本。

1. 这不是普通的利率波动，而是市场在重新定义“安全资产”的边界

传统上，主权债券被视为无风险资产的基准。但当财政可持续性受到质疑时，这一基准本身开始动摇。BIS论文的核心识别策略恰恰抓住了这一点：当财政风险上升时，投资者会从主权债券转向同等期限的最高等级公司债券。这种“组合再平衡”行为，使得主权债券收益率上升的同时，安全公司债券收益率下降——两者走势背离，正是财政风险冲击的独特指纹。

这一识别的精妙之处在于，它排除了货币政策、通胀预期或全球风险偏好等共同因素。如果主权债券收益率上升是因为央行加息，那么公司债券收益率也会同步上升；如果是因为通胀预期上升，两者同样会一起走高。只有在财政风险这个特定维度上，才会出现主权债券被“降级”、安全公司债券被“升级”的分化格局。

KC评论： 这意味着，投资者不再把本国国债当作终极安全资产。当市场开始用“公司债思维”来定价主权信用时，政府面临的融资约束将发生质变——不再是利率多高的问题，而是能否获得融资的问题。这对高债务国家是一个明确的预警信号。

BIS的实证结果描绘了一幅清晰的滞胀图景。财政风险冲击发生后，通胀和通胀预期立即上升，工业产出在短期内短暂增加，但随后进入持续性下滑。主权收益率曲线陡峭化，汇率贬值，股票价格下跌。

这一传导链条的逻辑是：财政风险冲击往往伴随财政扩张，短期内确实能刺激需求、推高通胀。但问题在于，主权债券收益率曲线整体上移（不只是短端），意味着整个经济体的融资成本都在上升。企业借贷成本增加、投资意愿下降，而汇率贬值又进一步加剧输入性通胀。结果是，通胀压力持续的时间远长于产出增长的时间——约六个季度后，更高的收益率和借贷成本导致实际活动持续放缓，股票估值下降。

论文特别指出，这种效应在财政风险冲击为“正向”（即财政状况意外恶化）时，比“负向”（财政状况改善）时更强且更持久。这意味着，财政扩张带来的通胀效应具有不对称性：恶化时通胀上升的幅度，远大于改善时通胀下降的幅度。这为“财政扩张易放难收”提供了新的量化证据。

论文最具有政策冲击力的发现，在于货币政策态度的调节作用。当央行在财政风险冲击后维持宽松立场（定义为政策利率调整幅度低于泰勒规则隐含的应有水平）时，滞胀效应被显著放大。

具体机制是：宽松货币政策在短期内确实能部分对冲借贷成本上升，但代价是实际利率持续为负。负实际利率进一步助长了通胀预期，同时未能阻止收益率曲线长端的上行——因为市场开始担心，宽松的货币环境会鼓励政府更无节制地扩张财政。约一年后，所有期限的收益率都显著上升，金融条件反而比不放松时更加收紧。

KC评论： 这给当前全球央行出了一道难题。当财政风险上升时，如果央行选择加息抗通胀，可能加剧经济下行；如果选择维持宽松，则可能陷入“财政主导”的陷阱，让通胀预期脱锚。BIS的结论指向一个反直觉的方向：在财政风险冲击面前，央行越“配合”，结果越糟糕。

论文还发现，财政风险冲击的效应高度依赖于初始条件。在主权C

[... middle omitted ...]

始主权风险溢价越高，财政风险冲击的破坏力越大。这类似于金融领域的“脆弱性叠加”：在低风险状态下，冲击可能被吸收；但在高风险状态下，同样的冲击可能触发系统性重定价。这对新兴市场国家尤其具有警示意义。

尽管BIS论文在识别和量化方面做出了重要贡献，但仍有几个关键问题留待进一步研究。

第一，财政风险冲击的“结构性来源”尚未被充分拆解。不同来源的财政风险（例如，是源于支出扩张、税收削减，还是隐性担保的显性化）可能产生截然不同的宏观效应。论文的识别框架虽然干净，但无法区分财政风险冲击的“成分”。

第二，公司债券的传导渠道仍是一个“黑箱”。论文假设投资者从主权债券转向安全公司债券，但公司债券市场本身的结构性特征（如流动性分层、信用评级迁移、投资者结构）如何影响这一传导，尚未被充分探讨。

第三，论文覆盖的12个经济体主要是发达经济体，新兴市场国家的财政风险传导机制可能完全不同——主权评级本身较低、公司债券市场深度不足、资本流动的脆弱性更高，这些因素都可能改变冲击的传导路径。

\subsection{[R012] 国际清算银行：高债务正在悄悄削弱加息的通胀抑制力}
{\small\color{refgray}归类：宏观与利率：亚洲央行轮动与全球财政风险}

市场主流叙事一直把通胀能否回落押注在央行加息力度上。但一份来自国际清算银行的最新工作论文给出了一个令人不安的答案：加息对通胀的抑制作用，可能被一个被长期忽略的结构性变量悄悄削弱——公共债务的规模和期限结构。

这份编号1365的工作论文，基于2001至2020年欧元区高频货币政策冲击数据，使用面板局部投影方法，得出了一个对全球投资者和政策制定者都极具冲击力的核心判断：高债务经济体面对加息时，价格和通胀预期的反应显著弱于低债务经济体，而产出收缩幅度却至少持平。换句话说，同样的加息幅度，在高债务国家，通胀更难压下去，但经济代价并不更小。

这不是一个学术象牙塔里的清谈。对于正在评估欧洲资产定价、全球利率路径和主权信用风险的决策者而言，这份报告提供了一个必须纳入分析框架的变量——债务期限结构如何改变了货币政策的传导效率。

报告的核心实证发现可以用一句话概括：当公共债务占GDP比例从60\%上升到120\%时，同样幅度的加息，价格水平的下降幅度显著缩小，通胀预期的锚定能力也明显减弱。

这一结果在统计上是显著的。报告作者使用了一个巧妙的识别策略：在保持债务期限结构不变的前提下，仅改变债务总量水平，来分离“债务规模”和“债务期限”两个变量的独立影响。结果显示，高债务经济体面对加息时，价格反应更弱，而产出收缩幅度至少与低债务经济体相当。

这意味着什么？传统教科书模型假设货币政策通过利率渠道影响总需求，进而传导至价格。但高债务环境下，这个链条被扭曲了。加息不仅抑制私人部门需求，还会通过影响政府资产负债表、债券持有人收入和财政风险溢价，产生一系列反向作用力。

KC评论： 简单理解就是，在高债务国家，央行加息就像在泥潭里踩油门——车轮在转，但车子没怎么动。通胀压不住，但经济的痛苦一点不少。这份报告用 年欧元区数据验证了这一机制，但更值得关注的是，它提供了一个分析框架：你可以用这个框架去评估当前美国、日本、意大利等国的货币政策有效性。

如果债务总量只是第一层变量，那么债务期限结构就是第二层更精密的调节阀。报告将政府债务按剩余期限分为四个桶：9个月以内、9个月至4年、4至8年、8年以上。结果发现，不同期限的债务对货币政策传导的影响呈现出高度非线性特征。

最关键的发现是：处于中间期限（4至8年）的债务，与最弱的产出和价格响应相关；而极短期（9个月以内）和极长期（8年以上）的债务，则与最强的响应相关。

这一非线性结果背后是三种机制的博弈。极短期债务使财政成本和债券持有人收入对利率变化高度敏感，强化了“利息收入效应”；极长期债务使债券价格对收益率变化更敏感，强化了“估值效应”；而中间期限的债务，可能使这两种效应部分相互抵消，从而削弱了货币政策的整体传导效率。

这一发现对资产定价有直接含义。如果某个主权国家的债务期限结构集中在中间区间，那么该国央行的加息对通胀的抑制作用可能比预期更弱，但经济衰退的风险未必更低。投资者在定价主权信用风险和汇率时，需要把这个维度纳入考量。

KC评论： 报告没有完全展开的一个问题是，不同期限债务的持有人结构差异如何影响这些机制。比如，如果中间期限债务主要由国内银行持有，而极长期债务主要由养老金和保险公司持有，那么估值效应和收入效应的传导路径会完全不同。这部分分析在完整报告中值得仔细阅读，它可能改变你对特定国家风险敞口的判断。

报告对“财政支持”的检验是另一个值得关注的发现。如果加息后政府能够同步收紧财政（即提供“财政支持”），那么债务对货币政策传导的干扰可能被抵消。但实证结果恰恰相反：加息后，政府的基本财政余额反而恶化，债务占GDP比例上升。

这意味着，在高债务环境下，央行加息不仅没有获得财政的配合，反而被财政的被动恶化进一步放大了政策困境。报告将其称为“缺乏同期财政支持”。这一发现与Sims（2

[... middle omitted ...]

次检验了欧元区货币政策对这些国家的溢出效应如何受到接收国债务期限结构的影响。

结果显示，溢出效应不是均匀的。当接收国债务中4至8年期限的占比更高时，欧元区加息对其产出和价格的收缩效应更强。这验证了估值效应在国际传导中的重要性：欧元区加息导致接收国长期债券价格下跌，通过估值损失渠道进一步收紧金融条件。

这一发现对全球投资者有直接意义。如果你在评估一个非欧元区欧洲国家的资产风险，不仅要看该国自身央行的政策，还要看其公共债务的期限结构——因为它决定了欧元区政策冲击的放大倍数。

KC评论： 报告构建的这个新数据集本身就是一个有价值的资产。它覆盖了保加利亚、克罗地亚、捷克、匈牙利、波兰等10国，时间跨度从1999年到2020年。对于研究新兴欧洲市场的投资者来说，这个数据集可能比很多公开数据更精细。完整报告附录中有详细的数据来源和构建方法，值得研究者直接使用。

尽管这份报告在实证设计上非常严谨，但它也留下了几个值得追问的问题，这些恰恰是读者在阅读完整报告时需要特别关注的。

\subsection{[R013] 波士顿咨询：CEO的AI红利，不在效率，在判断}
{\small\color{refgray}归类：企业AI转型：CEO亲自下场与组织重塑}

当一家公司谈论AI转型时，最值得观察的信号不是它的技术路线图，而是CEO本人如何使用AI。

波士顿咨询（BCG）在2026年6月发布的最新报告中，提出了一个被多数企业忽视的判断：AI对企业最高决策层的最大价值，不是节省时间，而是放大CEO及其核心团队的判断力。报告指出，尽管72\%的CEO现在直接负责公司的AI决策，但只有15\%的人从中获得了有意义的回报。这15\%的先行者与落后者的关键区别，不在于他们是否部署了更先进的系统，而在于他们每周至少花8小时亲自使用AI、理解其能力边界，并将其嵌入自己的决策流程。

这个发现颠覆了常见的AI转型叙事。多数企业将AI视为中层和基层的效率工具，用来优化流程、提升生产力。但BCG认为，真正未被开采的金矿在组织顶端——那里做出的决策影响公司未来数年走向，而判断质量的微小提升，都能带来几乎无限的商业价值。

1. 先行者CEO已经用AI扮演四种角色，但多数企业仍在观望

BCG通过研究15位高知名度CEO的公开案例，提炼出先行者使用AI的六种常见行为：快速掌握陌生领域知识、预判关键对话、压缩复杂信息、压力测试自身假设、更清晰地表达观点、以及评估自己的时间分配是否匹配战略优先级。

这些行为可以被归纳为四种角色：AI充当“随需导师”，帮助CEO在重大会议前快速建立对新领域的语感；充当“魔鬼代言人”，在决策尚未固化前挑战假设、揭示风险；充当“战略编辑”，将粗糙的想法转化为清晰的沟通；以及充当“反思工具”，通过分析日程和邮件来检验CEO的时间是否真的花在了最重要的事情上。

KC评论： 这些案例听起来像“高级秘书”或“个人助理”的替代品，但重点在于“压力测试”和“反思工具”。当CEO需要下属或董事会提出不同意见时，往往面临人际障碍——没人愿意当面挑战老板。而AI可以毫无负担地扮演这个角色。报告指出，这正是AI在CEO层面最被低估的价值之一：它能够提供人类顾问出于礼貌、偏见或信息距离而不愿提出的风险提示。完整报告中详细拆解了这些案例的具体操作流程，以及先行者CEO如何设计提示词来获得真正有价值的“反对意见”。

但BCG警告，当前多数CEO的AI使用仍然是“即兴的、用胶带拼凑的”，而非系统化的企业级部署。这意味着，先行者与追赶者之间的差距，正在从“是否使用”转向“如何使用”。

报告的核心洞察之一，是AI在CEO层面的应用正从通用工具向定制化Agent系统演进。BCG描绘了一个正在发生的范式转变：传统CEO做出重大战略决策的流程，通常始于一堆零散的预读材料、相互矛盾的备忘录和不一致的假设，随后是漫长的会议，大部分时间花在“对齐信息”上。

定制化AI Agent系统能够从根本上改变这个流程。通过训练在CEO个人的优先事项、决策历史、思维模型和持续变化的战略背景上，这些系统可以在会议开始前完成信息整合、冲突识别、权衡分析和风险预警。CEO走进会议室时，不再是“准备被汇报”，而是“准备做决策”。

KC评论： 这个“决策前准备室”的概念非常关键。它意味着AI的价值不是替代CEO的判断，而是大幅提升判断的质量和速度。BCG举了一个EMESA地区工业公司的案例，该公司正在构建一套定制化Agent系统，旨在为CEO提供实时绩效报告、风险分析、竞争情报和异常预警。当这套系统全面运行时，CEO将拥有组织的“实时脉搏”，能够提前发现潜在危机并将其转化为战略优势。完整报告中包含了这个案例的技术架构细节和初步成效数据，值得关注的是，这种系统如何平衡“信息全面性”与“认知过载”之间的矛盾。

然而，BCG也明确划出了边界：AI可以辅助决策，但不能为决策负责。责任始终在CEO身上。这也引出了下一个关键问题——当CEO越来越依赖AI时，如何避免被误导。

BCG用一整节篇幅警告了AI可能误导领导者的四种方式，每一

[... middle omitted ...]

相同的工具、相同的数据和相同的提示词，会强化相同的假设。BCG的研究发现，生成式AI的通用化输出会使团队的思维多样性降低41\%。这意味着一群聪明人可能更快、更一致地走向一个错误的共识，而人类的不同意见变得更加不可或缺。

第四种风险是“制造认知过载”。AI能够生成更多信息、观点和合成见解，但人的认知容量是有限的。BCG对美国1488名全职工作者的研究发现，14\%的AI用户报告了“AI脑疲劳”现象——因过度使用或监督AI而超出认知能力导致的疲惫。

KC评论： 这四种风险中，“思维多样性降低41\%”这个数据尤其值得警惕。它揭示了一个反直觉的困境：当所有人都使用相同的AI工具时，组织的集体智慧可能不是提升，而是下降。CEO需要主动设计“异议机制”——比如让另一个模型来批判第一个答案，或者在核心团队中指定一个“挑战者”角色。完整报告对这四种风险都给出了具体的应对框架，但最核心的建议是：CEO不能成为对抗误导性AI的最后一道防线，整个高层团队都需要具备“质疑AI答案”的纪律。

\subsection{[R014] 波士顿咨询：金融科技重回增长，但赢家的门槛已经变了}
{\small\color{refgray}归类：地产与消费：金融科技与游戏行业的结构性复苏 / 金融科技与保险：盈利修复与市场分化}

金融科技行业在经历2023至2024年的估值修正与资本寒冬后，2025年交出了一份出人意料的成绩单。全球金融科技总收入突破5000亿美元，同比增长22\%，增速是传统金融机构的四倍以上。但这份波士顿咨询与FT Partners联合发布的《2026年全球金融科技报告》传递的核心信号，并非简单的“行业复苏”。真正值得关注的判断是：金融科技已经从一个靠烧钱换增长的赛道，进入了一个利润与规模必须同时兑现的新阶段。行业没有回到2021年的狂热，而是进入了一个更成熟、但竞争门槛更高的“选择性繁荣”时期。

这份报告之所以重要，是因为它系统性地回答了三个投资者最关心的问题：增长是否可持续？哪些细分领域真正跑出来了？以及，AI和数字资产到底是泡沫还是下一波基础设施？在行业情绪从悲观转向谨慎乐观的当下，波士顿咨询的判断为决策者提供了一个难得的全局视角。

报告中最具说服力的数据，不是总收入的增长，而是运营质量的改善。全球最大的85家上市金融科技公司，2025年平均EBITDA利润率提升了4个百分点，达到20\%。更关键的是，74\%的公司已经实现盈利，而2024年这一比例是68\%。同时，符合“40法则”（收入增速加利润率超过40\%）的公司占比也在提高。

这意味着什么？2025年的增长不是靠烧钱买来的，而是靠真实的运营效率提升。资本市场的态度也印证了这一点：虽然股权融资总额回升至580亿美元，同比增长53\%，但资金高度集中在后期轮次。2023至2025年间，E轮及以后的融资额增长了超过210\%，而种子轮和天使轮则收缩了约10\%。投资者不再为故事买单，他们只愿意为已经证明有规模化路径和清晰经济模型的公司加注。

KC评论： 金融科技行业正在经历一个残酷但必要的“优胜劣汰”。2021年那种只要沾上“数字”概念就能拿到大额融资的时代已经彻底结束。今天，投资者要求的是可验证的利润路径，而不是PPT上的TAM（潜在市场规模）。对于产业决策者而言，这意味着并购整合的窗口已经打开——那些仍在亏损但拥有技术或客户基础的尾部公司，正在成为头部玩家低成本获取能力的标的。

2. 交易与投资、存款是增速最快的两极，但支付仍是绝对主航道

报告将金融科技收入按垂直领域拆解，揭示了明显的分化。交易与投资板块以38\%的同比增速领跑，存款紧随其后，增速达30\%。这两个板块的爆发，与2025年全球数字资产市场的回暖密切相关——加密市值重回约3万亿美元，稳定币市值约3000亿美元，代币化现实世界资产也达到约300亿美元。

但报告也提醒，支付依然是金融科技收入的绝对主力。在年收入超过5亿美元的大型金融科技公司中，支付板块的占比遥遥领先。这意味着，虽然新兴赛道在快速增长，但行业的基本盘依然建立在最传统的支付基础设施之上。这背后是一个尚未被充分打开的B2B机会：金融科技目前仅渗透了全球银行与保险收入池的4\%，而在B2B垂直领域，渗透率更低。

KC评论： 交易与投资的高增长，很大程度上是数字资产市场情绪修复的映射。但报告也明确指出，目前65\%的稳定币持仓仍与加密交易挂钩，数字资产的应用场景依然高度集中。这提示我们，如果剔除交易和投机需求，数字资产在支付、汇款等实体经济场景中的落地，距离规模化还有相当距离。完整报告中对代币化资产和稳定币的深入分析，值得所有关注Web3金融的读者细读。

3. AI的落地正在从“可能性”转向“可验证”，但赢家可能不是所有人想象的那样

AI是这份报告中最受关注的主题之一。波士顿咨询的结论是：AI确实在重塑金融服务，但当前的价值兑现集中在运营和工作流密集型领域，而非面向消费者的全自动体验。在承保、反欺诈、反洗钱、KYC、文档提取、客服、软件开发以及合规和财务工作流中，AI已经带来了可测量的效率提升。报告特别提到，使用AI的金融科技公司，开发者生

[... middle omitted ...]



这份报告尚未完全回答的一个关键问题是：AI驱动的金融科技公司，是否会像上一代“数字原生”公司一样，在规模扩张后遭遇相似的合规和资本效率瓶颈？AI能否让金融科技在盈利性上真正超越传统银行？这需要更长时间的观察，也是加入社群讨论后可以进一步追踪的话题。

从区域来看，亚太地区以25\%的增速成为全球增长最快的金融科技市场，日本和韩国的数字银行与加密交易平台是主要驱动力，东南亚的新加坡和印尼也贡献显著。欧洲以24\%的增速紧随其后，受益于监管环境的改善和新型银行向相邻产品与地理区域的扩张。

但报告也指出，区域间的分化不仅取决于需求，更取决于监管。在美国，大型金融科技公司正在积极寻求通过货币监理署获得联邦直接监管，以绕过合作银行的中间成本，实现更快的产品创新。类似趋势在欧洲也已出现。而在中国、印度和海湾部分地区，严格的监管环境仍使金融科技规模化面临更大挑战。

这种监管分化，正在重塑全球金融科技的竞争格局。那些能够在更复杂监管环境中运营的公司，反而可能获得更持久的竞争优势。

\subsection{[R015] 波士顿咨询：美国P\&C保险的“好日子”结束了，但赢家还没定}
{\small\color{refgray}归类：金融科技与保险：盈利修复与市场分化}

美国财产与意外险（P\&C）市场正在经历一个关键的转折点。波士顿咨询（BCG）在最新发布的《2026保险价值创造者报告》中给出一个核心判断：支撑过去几年行业高回报的“硬市场”红利正在消退，费率走软与结构性风险上升正在重塑竞争格局。这不是一次周期性的小波动，而是一次市场底层逻辑的切换。

报告显示，2021至2025年间，美国P\&C保险公司实现了约18\%的年化总股东回报率（TSR），跑赢全球保险业平均水平的15\%，也跑赢了大多数其他行业。但进入2025年，这一势头明显放缓。更值得关注的是，行业内头部与尾部公司的差距正在拉大，这种分化在历史上往往预示着一次市场洗牌的开始。

对于产业决策者和投资者而言，当下最需要回答的问题不是“市场会怎么走”，而是“在费率走软、风险复杂化的新环境下，什么样的公司能持续创造价值”。BCG的报告给出了一个清晰的答案框架，但其中的细节和隐含判断，值得逐层拆解。

KC评论： 波士顿咨询这份报告的价值不在于告诉你市场在变——这已经是共识。它的真正贡献在于，用五年期的TSR数据和承保利润率拆解，量化了“赢家”和“输家”之间的差距究竟来自哪里。读完你会发现，承保能力才是唯一的护城河，投资收入只是止痛药。

BCG报告中最具冲击力的图表之一，是不同分位保险公司对有形股本回报率（RoTE）的贡献拆解。在2021至2025年间，前四分之一的P\&C公司，其承保业务对RoTE的贡献约为11个百分点；而尾部公司的承保贡献为负值。

这意味着什么？在硬市场周期中，行业整体看起来都在赚钱，但利润来源完全不同。头部公司靠的是承保纪律——更低的综合成本率、更精准的风险定价；而尾部公司很大程度上依赖投资收益来掩盖承保端的亏损。

报告特别指出，2024至2025年间高企的净投资收益部分掩盖了弱势公司的承保恶化。随着利率正常化、投资收入回落，承保质量将成为区分盈利能力的唯一标尺。那些在硬市场期间投资于定价精细化和风险分层的公司，在费率走软的环境中更有能力捍卫利润率。

这是一个残酷的筛选机制：当潮水退去，裸泳者将无所遁形。对于投资者而言，判断一家P\&C公司的未来表现，不应看其当期利润，而应穿透利润结构，看其承保端是否真正具备定价权和风控能力。

美国个人车险市场是这份报告中最生动的案例。2022年，行业综合成本率飙升至112\%；到2025年，这一数字已改善至92\%——这是自2020年以来的最佳表现，也是2004年以来首次低于95\%的年份。修复速度之快，令人侧目。

但修复的成果分配极不均衡。报告点名了Progressive：这家公司率先、也最果断地重新定价其保单组合，2025年个人车险综合成本率降至约88\%，并实现了显著的保费增长，目前已成为美国个人车险市场的第一名。而其他公司则经历了更长时间的修复期——放弃保单、恢复利润率——之后才将综合成本率降至90\%出头。

现在，随着盈利能力基本恢复，竞争态势正在转变。保险公司正从“防守”转向“进攻”，重新进入此前退出的州和细分市场，重新启动增长杠杆，并更积极地竞争新业务。但BCG同时提醒，这种转变伴随着几个风险：索赔频率正从疫情低点回升，损失严重程度因车辆复杂性、维修成本和医疗费用通胀而持续上升；行业精算师估计的年损失成本增长率为6\%至8\%，在一些州已超过费率上涨速度；大型互助保险公司的竞争压力依然显著。

KC评论： Progressive的案例说明一个道理：在周期转换中，先动手、下重手的公司能获得结构性的市场地位提升。那些“等等看”的公司，即使最终也修复了利润率，却失去了市场份额。这份报告里的Progressive数据曲线和图表演示了“先发优势”在保险业中的具体形态，完整报告里有更详细的各公司对比图表。

对于车险市场的下一步，关键在于：当所有人都开始争夺新业务时，定价纪

[... middle omitted ...]

性重置”。

推动这一结果的关键因素是大飓风未登陆美国本土。即便2025年仍有帕利塞兹和伊顿山火等重大灾害，以及持续的对流风暴（SCS）损失，但缺少一次飓风登陆足以让行业整体数据看起来光鲜。报告警告，随着费率充分性在各州恢复，费率上涨动能将减速；损失成本通胀仍然活跃；山火和其他非模型化灾害预计将继续对损失率施压。

真正值得关注的是SCS风险的变化。根据Verisk和慕尼黑再保险的数据，2023、2024、2025年每年由SCS事件导致的保险损失均超过500亿美元，而此前十年的年均水平约为300亿美元。这些损失在地理上是分散的，使得投资组合管理和再保险结构比传统的沿海飓风损失管理更为复杂。

这意味着，房屋险市场的未来将出现明显的分化：拥有卓越巨灾模型、严格的敞口管理和充足定价能力的保险公司将获得强劲回报，而其他公司的业绩将被未来的自然灾害事件侵蚀。规模正在成为护城河——前五大个人险公司的市场份额优势持续扩大，而小型公司在理赔处理、分销和技术上的固定成本劣势日益凸显。

\subsection{[R016] 波士顿咨询：95\%的公司正在错过AI真正的价值拐点}
{\small\color{refgray}归类：AI与算力：存储器超级周期与AI落地组织变革 / 企业AI转型：CEO亲自下场与组织重塑}

一份覆盖全球1,250家企业的实证研究给出了一个令人不安的判断：只有5\%的公司真正从AI中获得了可量化的商业价值。这5\%的企业——波士顿咨询称之为“未来型公司”——不仅跑赢了同行，而且正在加速拉开差距。它们的收入增速是落后者的1.7倍，三年股东总回报是后者的3.6倍，投入资本回报率高出2.7倍。

这份报告的核心洞察并不在于“AI有用”，而在于“AI的价值分配已经出现结构性断裂”。60\%的企业虽然持续投入，但几乎没有获得任何实质性回报。真正的问题不是技术不够好，而是大多数公司用错了策略框架——它们在做自动化，而领先者在做再造。

1. 5\%的公司正在构筑一条AI价值护城河，而60\%的公司正在被反向锁定

波士顿咨询的研究将企业按AI成熟度分为四类：未来型（5\%）、规模化中（35\%）、新兴型（46\%）和停滞型（14\%）。后三类合计95\%的企业，构成了“追赶者”阵营。

最关键的数字不是5\%这个比例，而是这些未来型公司正在创造一种“复合效应”。它们将AI产生的利润重新投入技术基础设施和人才，形成正向循环。2025年，这些公司计划将IT预算中高达64\%的份额投入AI，而追赶者这一比例不足其一半。结果是，未来型公司预期在AI应用领域实现双倍于追赶者的收入增长，成本降幅也高出1.4倍。

对于那60\%几乎没有回报的公司来说，问题出在顶层。波士顿咨询指出，这些公司的管理层往往“口头上重视AI，但没有明确的量化目标，也没有定期追踪进度的机制”。他们将AI决策下放到中层或基层，导致资源分散在数十个孤立的试点项目中，无法形成规模价值。

KC评论： 这份报告揭示了一个被低估的风险——AI投资正在变成一种“富者愈富”的游戏。未来型公司每赚1块钱的AI回报，就会再投1.26元去巩固优势。而落后者的困境在于：它们越不投入，就越看不到效果；越看不到效果，就越难说服管理层继续投入。这不是技术问题，是战略节奏问题。报告中对不同区域和行业的成熟度对比，以及具体的投资回报率数据，是理解这个循环的关键——完整版有更详细的行业拆解和区域差异图表。

2. AI价值高度集中在核心业务职能，但IT的占比正在意外跃升

报告给出了一个清晰的职能价值分布：70\%的AI价值潜力集中在销售与营销、制造、供应链、定价等核心业务职能。研发与创新单独贡献了15\%。这一趋势与2024年的发现一致。

但有一个意外的变化：IT职能的AI价值占比从2024年的7\%跃升至2025年的13\%，增加了6个百分点。波士顿咨询的解释是，随着AI从预测性向生成式和代理式演进，IT部门正在从“支持角色”转变为“价值创造中心”。尤其是代理式AI的部署，对数据基础、系统架构和治理能力提出了更高要求，这些恰恰是IT的核心能力域。

在具体工作流层面，报告给出了可量化的案例：能源公司通过AI驱动的预测性维护，在全面部署后可避免30\%的可控成本；一家全球美妆公司通过引入AI虚拟美容顾问，预期在12个月内获得1亿美元的增量收入，ROI是传统电商路径的两倍。

这些数字表明，AI价值不是均匀分布的，而是集中在少数几个经过精心选择的工作流中。未来型公司的做法是“端到端重塑”而非“点状优化”。

3. 代理式AI正在成为加速差距扩大的核心变量，但部署顺序决定成败

报告中最具前瞻性的部分是关于代理式AI的分析。这类AI结合了预测性和生成式能力，能够自主推理、学习和执行任务。2024年几乎无人提及的代理式AI，在2025年已占AI总价值的17\%，预计到2028年将升至29\%。

但波士顿咨询给出了一个重要的警示：代理式AI不是AI部署的起点，而是进阶阶段。其前提条件包括强大的数据基础、已规模化的AI能力、清晰的治理框架。最成功的案例是将代理嵌入整个工作流，而不是作为孤立的试点运行。

一家全球电子设备制造

[... middle omitted ...]

裁决，失败责任如何归属——这些是完整报告里值得细看的附录内容。

波士顿咨询明确指出，未来型公司的策略是可复制的，不是黑箱。它们遵循五个相互关联的策略：设定多年度战略AI愿景、重塑并创造端到端影响、采用AI优先的运营模式、保障并赋能必要人才、使用适配的技术与数据架构。

听起来不复杂。但问题出在“同时做到”。报告的数据显示，追赶者中最常见的失败模式是：管理层有愿景但没有量化目标，IT部门有工具但没有业务部门的协同，员工有AI工具但没有足够的培训。超过50\%的未来型公司计划在未来一年内对内部员工进行大规模技能提升，而落后行业的这一比例不到15\%。

区域差异也值得关注。亚太地区的公司在IT预算中分配给AI的比例最高（5.2\%），并预期到2028年实现10\%的收入增长和12\%的成本降低，均高于北美和欧洲。但在代理式AI的采用率上，北美领先（51\%），亚太次之（45\%），欧洲落后（41\%）。这暗示了一个有趣的动态：亚太在“投入意愿”上更激进，但在“先进工具采用”上仍有差距。

\subsection{[R017] 波士顿咨询：游戏行业的下一个增长引擎不是硬件，是平台碰撞}
{\small\color{refgray}归类：地产与消费：金融科技与游戏行业的结构性复苏}

游戏行业正在走出“三年寒冬”。这不是一个模糊的乐观预期，而是波士顿咨询（BCG）在最新发布的《Video Gaming Report 2026》中给出的明确判断。这份基于近3000名全球玩家调研和行业深度分析的报告，核心结论只有一个：游戏行业的下一轮增长，不会来自新一代主机或更快的芯片，而是来自四大趋势的“碰撞”——生成式AI、用户生成内容（UGC）、云游戏和开放应用商店。当这些趋势同时发生，传统的“主机战争”将让位于“生态战争”，而赢家将是那些重新定义分发、发现和变现逻辑的公司。

报告预测，全球游戏收入将从2025年的约2600亿美元增长至2030年的约3300亿美元，复合年增长率约为5\%。这个数字远低于2010年代行业规模翻倍的速度，但BCG认为，这是更健康、更可持续的增长模式。更重要的是，增长的结构正在发生根本性变化：云游戏收入将从2025年的约15亿美元飙升至2030年的约300亿美元，UGC创作者经济仅Fortnite和Roblox两个平台在2025年的分成就将超过15亿美元，而应用商店的开放将重塑移动游戏50\%全球收入的分发逻辑。

1. 云游戏终于不再是一个概念，2025年是进入主流的分水岭

过去十年，云游戏一直被贴上“未来技术”的标签，但BCG的调研数据表明，这个未来已经到来。全球范围内，60\%的玩家尝试过云游戏，其中80\%给出了正面体验反馈。这个数字远超行业预期。

更关键的是，云游戏正在改变“平台”的定义。传统的主机战争——PlayStation vs Xbox vs Nintendo——本质上是对客厅硬件的争夺。但云游戏让硬件变得无关紧要。玩家可以在手机、平板、PC、智能电视甚至汽车屏幕上无缝游玩3A大作。BCG报告明确写道：“主机战争将变得越来越无关紧要，竞争正在转向由全屏云游戏技术支撑的生态系统。”

这意味着什么？对于硬件厂商，未来五年的投资逻辑必须从“卖更多主机”转向“运营更好的云平台”。对于游戏开发者，不再需要为不同平台做大量适配工作，分发成本将大幅下降。对于投资者，那些在云基础设施、流媒体技术和跨平台中间件上布局的公司，将获得超越行业平均的回报。

KC评论： 云游戏的普及速度比大多数人想象的要快。80\%的正面体验率说明技术瓶颈正在被突破。但报告没有完全展开的是：云游戏的商业模型仍然不清晰——是订阅制、按小时计费还是广告支持？这个问题将直接影响平台方的利润率和竞争格局。完整报告中有详细的云游戏收入分拆和三种商业模型对比，值得细读。

UGC（用户生成内容）在过去几年被很多人视为“小孩子玩的东西”——Minecraft和Roblox的典型用户是8-12岁的儿童。但BCG的调研揭示了一个更重要的趋势：UGC正在渗透到所有年龄段。

报告数据显示，40\%的玩家表示他们消费的UGC内容比一年前更多。更重要的是，UGC正在成为新一代玩家的“第一入口”。大约44\%的玩家父母表示，他们的孩子在5岁之前就开始玩电子游戏，而孩子玩的最多的前三款游戏中，有两款是UGC游戏：Minecraft和Roblox。

这意味着UGC不仅是内容形式，更是一种用户获取渠道。当新一代玩家从小就在UGC生态中长大，他们对“什么是游戏”的定义将与上一代完全不同。他们不关心画质有多好、剧情有多复杂，他们关心的是能否创造、能否分享、能否与朋友一起玩。

KC评论： 这是一个代际转换的信号。如果UGC成为主流游戏体验，那么传统的“卖拷贝”模式将面临根本性挑战。未来的游戏公司可能更像内容平台——核心能力不是开发，而是运营创作者生态和分发网络。报告中有关于Fortnite和Roblox创作者经济的具体数字，以及UGC如何影响游戏设计原则的分析，这些细节在完整报告中才有。

3. 应用商店的开放是一场“地震”，但震中在移动

[... middle omitted ...]

所未有的机遇；对于大型发行商，这意味着需要重新思考市场推广预算的分配。

报告同时也指出，这种开放不会一蹴而就。监管的推进速度、用户支付习惯的改变、以及平台方的反制措施，都将影响最终格局。但方向是确定的：应用商店的围墙花园正在倒塌。

BCG报告对AI的讨论非常务实。它没有陷入“AI是否取代人类开发者”的争论，而是聚焦于一个更关键的问题：AI将如何改变游戏市场的供给结构。

报告引用Steam数据，截至2025年中期，约20\%的新游戏披露使用了AI，这个比例是一年前的两倍。BCG估计约50\%的工作室已经在使用AI工具。更值得关注的是，AI-native工作室正在崛起——例如BitMagic和Series Entertainment，它们声称可以加速游戏开发90\%。

但BCG同时发出了一个重要的警告：AI会催生大量低质量内容，报告称之为“gameslop”。当AI可以快速生成大量平庸游戏时，发现和策展将成为平台的核心竞争力。玩家不会缺少游戏，但会缺少值得花时间的游戏。

\subsection{[R018] 麦肯锡：美国250年竞争力未断，但真正考验的是“下一章”能否续写}
{\small\color{refgray}归类：美国竞争力与劳动力市场转型}

麦肯锡：美国250年竞争力未断，但真正考验的是“下一章”能否续写

美国建国250年之际，麦肯锡全球研究院发布了一份罕见的全景式报告，回顾美国经济竞争力的历史弧线，并前瞻下一阶段的挑战。这份报告的核心判断出人意料地清醒：美国目前仍是全球最具竞争力的经济体，但多个历史优势正在变成负债。真正的问题不是庆祝过去，而是能否再次找到资源、雄心、制度与政策的新对齐，以续写下一个章节。

这份报告的价值不在于罗列美国有多强，而在于它提供了一个少见的、长周期的分析框架，把AI、地缘政治、人口结构、财政健康、教育衰退、基础设施老化等看似分散的变量，统一到一个“竞争力演化”的叙事中。对于关注全球资产配置、产业趋势与中美竞争格局的读者，这是一份值得细读的底层框架。

1. 美国竞争力的“账面”依然强劲，但结构性的裂缝已经清晰可见

报告开篇即点明：美国以全球4\%的人口，创造了26\%的全球GDP，贡献了超过50\%的全球市值。在全球前100强企业中，美国占据了59席。过去几年，美国的生产率增速在主要经济体中领先，宣布的绿地外商直接投资流入也比疫情前翻了一番。

这些数字固然亮眼，但麦肯锡并不满足于复述成就。报告在第一章就引入了“竞争力”的三维定义：全球领先的企业、创新与技术领导力、个人经济机会与繁荣。前两个维度美国表现突出，但第三个维度——经济机会——已经出现明显裂痕。收入不平等、教育水平下滑、基础设施评级仅为C级，这些都被报告明确列为“需要反思的理由”。

KC评论： 这份报告最值得注意的不是美国有多强，而是它把“竞争力”定义为一个包含社会包容性的概念。如果只看GDP和市值，美国依然遥遥领先；但如果把教育、基建、财政可持续性纳入，画面就复杂得多。完整报告中有几十个横向对比图表，覆盖G7国家与中国，值得逐一对照。

报告把美国250年的经济史划分为四个章节：农业时代、工业时代、科学时代、数字时代。每一章都对应一次经济模式的根本性重构。美国之所以能持续领先，不是因为固守某种优势，而是因为它能够反复“破坏并重建”自己的经济基础。

这一分析框架的洞察力在于，它解释了为什么简单的“美国衰落论”或“美国例外论”都过于粗糙。真正的竞争不是看某一年GDP增速的对比，而是看一个国家能否在技术范式转换时，快速重组制度、人力资本和基础设施。麦肯锡指出，美国在每一次转型中都依靠了两个深层基础：创新文化与自然禀赋。但这两个基础本身不会自动保证成功，它们需要被“集体行动”激活——也就是报告反复强调的“We the people”。

KC评论： 这段历史分期是报告最精彩的部分之一。它提供了一个理解“为什么美国能反复反弹”的框架，同时也暗示了下一个章节可能需要的条件。完整报告中有专门章节拆解每一次转型的关键驱动因素和政策选择，对于思考当前AI转型的投资者来说，非常有参考价值。

报告在第三章和第四章详细拆解了美国当前面临的五个核心挑战：财政健康恶化、基础设施老化、教育成绩下滑、制造业专长流失、收入与财富持续不平等。

以财政为例，报告指出，美国国债占GDP的比例已超过100\%，且利息支出正在挤占其他公共投资。基础设施方面，美国土木工程师学会给出的评级是C级，这意味着未来十年需要约7.2万亿美元的投资。教育方面，PISA成绩在发达国家中排名中游，且过去十年没有明显改善。制造业方面，尽管美国仍保留超过50\%的全球半导体收入份额，但制造产能已从1990年的37\%降至2022年的10\%。

这些挑战并非新问题，但麦肯锡的独特之处在于，它把这些“存量问题”与“增量变量”联系在一起：AI正在打开新的可能性空间，地缘政治竞争正在加剧，生育率正在下降。换句话说，美国不仅要解决旧账，还要同时应对新环境。

KC评论： 报告没有给出简单的“政策建议清单”，而是提醒读者：这些挑战是相

[... middle omitted ...]

尚未回答的子问题：第一，AI的产业化能力是否足够快，以弥补制造业外流带来的就业和技能缺口？第二，美国的教育体系能否培养出足够多的AI时代人才，还是将继续依赖移民？第三，财政和地缘政治压力会不会压缩政府和企业在AI基础设施上的投资空间？

这些问题的答案，将决定美国能否像前四次转型一样，成功进入下一个增长周期。报告没有提供确定答案，但提供了一个清晰的“追问框架”。

基于这份报告，我们可以提炼出一个更实用的观察框架。对于关注美国经济与全球市场的读者，不必过度纠结于季度GDP或通胀数据，而应该关注以下几个长期信号：

第一，美国能否在AI领域保持“从研究到产业化的完整闭环”？这需要关注研发投入结构、人才流动、以及联邦政府对AI基础设施的支持力度。第二，美国能否在基础设施和制造业回流上取得实质性进展，而不仅仅是宣布投资计划。第三，教育体系能否在下一个十年内出现明显改善，尤其是STEM领域的成绩和技能培训。第四，财政可持续性是否会被纳入政策优先级，还是继续被短期政治博弈拖延。

\subsection{[R019] 麦肯锡：HR正站在“被IT吞噬”与“重塑组织”的分叉口}
{\small\color{refgray}归类：美国竞争力与劳动力市场转型}

2026年，全球劳动力市场正在经历一个微妙的转折点。宏观经济的不确定性让员工流动率下降，AI的渗透让技能需求结构发生根本性位移，而HR部门自身的变革速度却远远落后于外部环境的变化。

麦肯锡最新发布的《HR Monitor 2026》报告，基于对全球10个国家、1300名HR专业人士和5500名员工的调研，给出了一个清晰的判断：HR职能正站在一个分叉路口。一条路是继续维持现状，让日益复杂的技术变革超出自身能力边界，最终导致HR的核心职责被IT或数字化部门逐步蚕食。另一条路则是主动重塑，成为AI时代组织转型的架构师和领航员。

这份报告最值得关注的判断并非关于某个具体的HR流程改进，而是一个系统性的警示：如果HR不能在战略规划、技能管理、运营模式这三个维度上完成从“职能卓越”到“系统级变革”的跃迁，它在未来三到五年内面临的根本不是地位提升问题，而是职能边界被重新定义的问题。

KC评论： 麦肯锡的这份报告本质上是在说，HR的“舒适区”正在消失。过去几年，HR在流程优化、招聘效率、员工体验等方面取得了不少进步，但这些进步大多是在既定框架内的改良。现在，AI带来的任务重组和组织形态变化，要求HR必须跳出“管人”的旧角色，去思考“人和智能体如何协作”的新命题。报告中的关键数据——比如只有11\%的企业在做三年以上的战略性人力规划——暴露了多数组织在这方面的准备严重不足。

1. 战略规划的最大盲区：11\%的企业在做长期规划，而技能缺口正在加速扩大

麦肯锡的报告揭示了一个核心矛盾：几乎所有企业都在做人力规划，但绝大多数是短期的、运营性的。数据显示，62\%的受访企业会进行组织层面的人力规划，但只有11\%的企业将规划视野延伸至三年以上，而近三分之二的企业规划周期不超过12个月。

这种短视在AI加速渗透的背景下尤为危险。根据麦肯锡全球研究院的分析，如果AI以中等速度被采用，到2030年，欧美30\%的当前工时可能被自动化。与此同时，超过70\%的当前技能预计仍将保持相关，但这并不意味着岗位本身不会发生变化。几乎每个职业都将经历技能结构的重组。

报告中的一组数据值得注意：HR专业人士认为23\%的员工缺乏当前岗位所需的技能，这一比例同比下降了9个百分点。但22\%的员工自己表示，怀疑自己在未来五年内能否具备留在当前岗位的技能。这两个数据都在下降，但方向值得警惕——HR似乎比员工更乐观。更关键的是，HR专业人士对未来技能变化的程度可能仍然低估了。

KC评论： 这里有一个容易被忽视的“乐观偏差”。当HR认为技能缺口在缩小，而员工自己的信心也在回升时，很容易让人觉得“问题没那么严重”。但麦肯锡的报告暗示，这种乐观可能是基于对AI影响速度的低估。真正的风险不在于当前缺口的大小，而在于未来技能需求变化的节奏。如果规划周期只有12个月，而技能结构在3年内可能发生30\%的重组，那么大多数企业实际上是在“边开车边换轮胎”。

2. 技能需求正在发生结构性位移：软件开发者不再是香饽饽，数字素养和推理能力在崛起

报告对未来技能需求的排序变化提供了一个极具洞察力的窗口。2026年的调查显示，问题解决能力仍然是最受重视的未来技能，44\%的HR将其列入前五。创造力从去年的第五位跃升至第二位。数据分析和AI继续稳居前三。

但真正值得关注的是那些排名大幅变化的技能。数字素养从2025年的第20位飙升至今年的第6位，成为上升最快的技能。推理能力从第12位升至第9位。与此同时，软件开发和其他高度专业化的技术执行技能的重要性在下降。

这个变化背后的逻辑很清晰：AI系统正在越来越多地承担具体的执行任务，包括编写代码。因此，人类需要的能力正在从“如何做”转向“如何用、如何判断”。数字素养不再是IT部门的专属技能，而成为所有知识工作者的基础能力。推理能力和创造力之所

[... middle omitted ...]

比下降了2个百分点，即便员工满意度整体保持稳定。

这听起来是HR的好消息——招人更容易，留人压力更小。但报告同时指出，招聘效率问题并未因此缓解。企业面临的一个新矛盾是：虽然候选人的竞争烈度下降，但招聘流程的复杂性和周期长度并未改善。更长的招聘周期意味着，即便市场偏向雇主，那些最优秀的候选人仍然可能因为等待时间过长而流失。

AI在招聘领域的应用潜力巨大，但报告发现，多数企业仍然停留在将AI“叠加”到现有复杂流程上，而非用AI重新设计流程本身。比如，AI可以大幅提高简历筛选和面试安排的效率，但如果企业的招聘流程本身就是一个多环节、多审批、多轮面试的“长链条”，AI只能加速其中某个环节，而无法解决整体效率问题。

KC评论： 这其实是一个很典型的“数字化悖论”。企业在引入AI时，往往选择在现有流程上做“加法”，而不是做“减法”。结果就是，流程变得更复杂而非更简单。麦肯锡报告隐含的提醒是：AI的价值不在于让一个糟糕的流程跑得更快，而在于让你有机会重新思考这个流程是否必要。

\subsection{[R020] 麦肯锡：2025年，企业最大的风险不是AI，而是把学习当成福利}
{\small\color{refgray}归类：未被主题摘要引用}

麦肯锡：2025年，企业最大的风险不是AI，而是把学习当成福利

2025年，企业面临的最大挑战不是技术选型，不是成本控制，而是如何让组织在持续动荡中保持“学习能力”这一核心资产。麦肯锡研究与创新学习实验室最新发布的《2025学习趋势展望》给出了一个反直觉的判断：当大多数企业还在纠结“要不要用AI替代员工”时，真正领先的组织已经在问一个更根本的问题——如何让工作本身变成发展的引擎。

这份报告不是一份趋势清单。麦肯锡明确将其定义为“导航工具”，而非“趋势摘要”。报告的核心主张是：未来三年，企业的人才发展必须从“支持功能”升级为“战略驱动力”，而实现这一转变的钥匙，藏在三个相互强化的趋势中——流动的发展生态系统、负责任的AI采纳、以及组织韧性的系统性构建。

报告中最值得关注的信号不是某个新技术，而是一个认知转折：发展不再是一项福利，而是一种“关怀行为”（act of care）。这个表述听起来温和，但其商业含义极为锋利——那些把员工发展当作福利来做的企业，将在人才争夺战中系统性落后。

KC评论： 这份报告最值得反复咀嚼的不是它列出的趋势，而是它提出的底层逻辑。麦肯锡明确说，三个宏观趋势——流动生态、负责任的AI、韧性——不是独立话题，而是“相互强化、无法孤立成功”的系统。这意味着，过去那种“今年主推AI培训、明年主抓员工幸福感”的碎片化策略，已经行不通了。完整报告里有一张非常清晰的系统互动图，值得仔细研究。

1. 流动发展生态：从“碎片化学习”到“工作即发展”的范式转换

报告的第一大趋势直指一个长期痛点：企业的人才发展职能仍然处于“被动支撑”状态。2023年ATD调查显示，虽然72\%的受访者认识到将HR转型为跨职能学科的重要性，但只有11\%报告了有意义的进展。这个数字本身就说明问题——认知与执行之间存在巨大鸿沟。

麦肯锡认为，未来的发展方向是“流动的发展生态系统”（fluid development ecosystems）。这不是一个模糊的概念，而是三个具体动作的组合：去孤岛化的人才职能、数据驱动的发展生态、以及基于前瞻性洞察的战略决策。

最值得关注的是“去孤岛化”这一条。报告指出，尽管L\&D、HR、人才等职能之间的协作意愿在增强，但“结构性和决策权仍在强化旧的边界。围栏没有倒下，只是移动了位置。”这意味着，真正的变革不是组织架构调整，而是权力结构和决策逻辑的重塑。

数据方面，报告引用了世界经济论坛《未来就业报告》的数据：2025年至2030年间，39\%的现有技能将被转化或过时。但更关键的是，61\%的企业仍然只做一年期的人力规划。这个对比揭示了企业面临的真实困境——明知技能在快速贬值，却缺乏改变的战略耐心。

KC评论： 这里有一个容易被忽略的细节。报告提到“技能验证”正在从自我报告数据转向更完整的验证体系，但没有展开。完整报告里有一张关于“技能数据生态”的图，展示了从数据采集到预测建模的完整链路。对于正在搭建人才数据体系的企业来说，这张图本身就是一套很好的架构参考。

第二大趋势“负责任的AI采纳”可能是这份报告中最具紧迫感的部分。麦肯锡的判断是：“信任是脆弱的，处理不当可能会侵蚀员工的信心。”这个表述背后有一个被很多企业忽视的现实——员工对AI的焦虑正在上升，而企业对此的沟通和引导严重不足。

报告指出，尽管人们领导者正在积极实验AI——构建AI Agent、建立合作伙伴关系、试点新功能——但这些努力“可能给员工带来更多复杂性，而不是明确的价值”。这是一个尖锐的批评：企业AI战略最大的风险不是技术失败，而是信任破裂。

麦肯锡给出的方向是：负责任地采纳AI，必须向员工传递一个明确信号——领导者关心他们的贡献，并致力于确保AI帮助他们成功。这听起来像是常识，但实践中极少有企业做到了。大多数AI转型

[... middle omitted ...]

次：个体韧性和组织韧性。个体韧性强调的是“恢复能力”，而组织韧性是“预期挑战、适应并成长”的能力。这个区分非常关键——前者是被动的，后者是主动的。

数据方面，报告引用了一个令人警醒的数字：2025年，66\%的员工报告工作倦怠。这个数字来自2025年2月的最新研究。在员工倦怠率持续攀升的背景下，仅仅提供“减压课程”已经不够。企业需要重新设计工作流本身——包括如何分配任务、如何设置目标、如何提供恢复时间。

报告特别提到“支持恢复以确保持续绩效”，并强调这需要“可持续的工作流”。这意味着，韧性建设不是HR部门的专项活动，而是一场涉及组织设计、流程再造和文化重塑的系统工程。

尽管这份报告提供了极具洞察力的框架，但有两个关键问题没有被充分回答。

第一，成本与回报的量化。报告提到“发展是最大的关怀行为”，但没有给出企业投入产出比的测算。对于CFO来说，“关怀”不是预算科目。当企业面临成本压力时，这类投入往往首当其冲。报告没有提供一个可量化的商业案例，这可能会限制其说服力。

\subsection{[R021] 麦肯锡：88\%企业在试AI，但81\%未见到利润}
{\small\color{refgray}归类：AI与算力：存储器超级周期与AI落地组织变革 / 企业AI转型：CEO亲自下场与组织重塑}

麦肯锡在2026年2月发布的《组织状态》报告中，抛出了一个让企业决策者无法回避的核心矛盾：绝大多数组织已经开始拥抱人工智能，但绝大多数组织也并未因此获得可观的财务回报。88\%的受访组织正在以某种形式试验AI，然而81\%的受访者坦言，这些尝试并未对利润产生实质性的影响。这不是一个关于技术是否有效的疑问，而是一个关于组织是否准备好迎接技术的问题。

这份报告基于对全球超过10,000名高管和领导者的调研，覆盖15个国家与16个行业。它试图回答一个根本性的追问：在技术、经济和劳动力三大结构性力量同时重塑商业环境的今天，组织如何才能从短期的韧性转向持续的生产力与长期的价值创造？

报告得出的结论清晰且尖锐：那些只将AI视为效率工具、在现有流程上做局部修补的组织，注定与真正的价值增长无缘。真正的赢家，将是那些敢于推动“双重转型”——同时改造技术与组织——的企业。它们需要重新想象工作如何完成，重新定义职能与端到端流程，并从根本上重塑传统结构。

KC评论： 这个数据点值得所有管理者停下来思考。88\%的试验率与81\%的无利润结果之间，存在一个巨大的“价值黑洞”。报告暗示，这个黑洞并非技术本身，而是组织层面的障碍——缺乏清晰的所有权、缺乏系统性的流程重构、以及缺乏对员工能力的同步投入。这意味着，企业现在面临的不是“要不要用AI”的选择，而是“如何用组织变革来兑现AI承诺”的紧迫命题。

麦肯锡的调研揭示了一个令人不安的认知鸿沟：高达86\%的领导者认为，自己的组织在将AI融入日常运营方面准备不足。更值得警惕的是，每六家组织中就有一家没有明确的C级负责人来推动AI采纳。这意味着，在许多企业中，AI战略要么沦为空谈，要么被降级为某个部门的技术实验，缺乏自上而下的系统推动。

报告指出，当前大多数AI整合努力都集中在碎片化的用例上，旨在提升个体贡献者的效率。那些试图将AI智能体嵌入现有流程以驱动生产力的更实质性努力，要么仍处于规划阶段，要么正在试点项目中测试。这种“头痛医头”的路径，恰恰是导致81\%的企业无法看到利润增长的根本原因。

麦肯锡的研究表明，那些真正成功的组织，其AI转型并非始于技术采购，而是始于对工作流程的端到端重构。它们不满足于用AI加速现有步骤，而是重新思考哪些步骤应该存在、哪些可以合并、哪些可以由AI智能体自主完成。这种系统性思维，才是从效率提升跨越到价值创造的关键。

KC评论： 86\%这个数字本身就是一个信号。它意味着大多数组织在AI转型的起跑线上就已经落后。而“没有C级负责人”这一发现尤其关键——它说明在许多企业里，AI仍然被视为IT部门的职能，而非CEO或COO的战略议程。报告后续章节对“人机协作”和“共享服务中心重构”的深入分析，正是为了回答“组织准备度”这个核心问题。这些内容在完整报告中配有详细的案例拆解和操作框架。

2. AI智能体不是工具，而是“同事”——但只有四分之一的领导者这么认为

麦肯锡的报告将AI智能体的角色定位提升到了一个全新的层次。它指出，AI智能体不仅仅是更聪明的工具，它们应该成为人类员工的“自主队友”——能够被插入公司工作流程，自主完成任务，甚至在多智能体协同的场景下，它们可以组成团队并设计自己的工作流。

然而，受访领导者的认知显然滞后于这一愿景。只有四分之一的领导者预期，在短期内AI智能体将作为员工的自主队友发挥作用。这种认知落差，直接阻碍了组织对AI潜力的释放。

报告引用了一个来自电信公司的案例：该公司构建了一个“下一个最佳体验”引擎。AI模型识别客户何时可能需要帮助或更优的报价，然后通过客户偏好的渠道发送个性化信息。只有当需要时，才会触发人工介入。这一系统降低了客户流失率，改善了利润率，并显著提升了客户参与度。这个案例的关键在于，AI并非替代了人类，而是重新定义了人类介

[... middle omitted ...]

来说，是绕不开的功课。

一个被大多数企业忽视的AI应用沃土，是共享服务中心。麦肯锡的报告指出，随着技术重塑工作方式，传统的共享服务中心正在从交易处理中心演变为全球业务服务中心。84\%的领导者计划在未来一到两年内扩大其共享服务中心的范围。然而，超过40\%的受访者尚未开始系统性地采纳所需的技术。

报告描绘的愿景是：这些中心将“AI优先”地设计，从物理中心变为虚拟中心，在人类与AI智能体之间编排工作，解锁端到端的自动化，并在规模上驱动创新和洞察。这意味着，财务、人力资源、IT支持等后台职能，将不再仅仅是成本优化的对象，而可能成为整个企业AI转型的试验场和动力源。

KC评论： 84\%的计划扩张率与40\%以上的技术采用滞后率之间，存在一个明显的“执行缺口”。报告暗示，那些率先完成共享服务中心AI化改造的企业，不仅能够获得成本优势，更可能建立起竞争对手难以复制的运营敏捷性。完整报告对共享服务中心的转型路径、技术选择和组织设计有更详细的拆解，包括具体的企业案例和投资回报测算。

\subsection{[R022] 麦肯锡：CEO不下场，AI永远只是成本}
{\small\color{refgray}归类：AI与算力：存储器超级周期与AI落地组织变革 / 企业AI转型：CEO亲自下场与组织重塑}

这份报告最值得看的判断不是“AI使用率又创新高”，而是：真正从AI中拿到利润的公司，正在做三件大多数企业还没做的事——把CEO拉进治理、重新设计工作流、以及用KPI而不是讲故事来管理AI项目。

麦肯锡2025年3月发布的《The State of AI》全球调查，覆盖1491位来自101个国家的企业受访者。报告试图回答一个核心问题：当几乎所有公司都在用AI时，为什么只有极少数看到了真实的利润影响？

答案不是技术选型，不是算力投入，而是组织架构的“重布线”。这份报告提供了迄今为止最系统的证据，说明企业需要做出哪些结构性的改变，才能把AI从“实验工具”变成“利润引擎”。

以下是我们从这份报告中提炼出的六个核心洞察，以及一个报告没有完全回答的关键追问。

报告最反常识的发现是：年收入超过5亿美元的大型企业，在AI部署上正在加速拉开与中小企业的差距。这不是简单的“大公司更有钱”，而是大公司在治理结构、风险管理和人才配置上采取了完全不同的路径。

大型企业更倾向于让CEO亲自监督AI治理，而中小型企业更多把这事交给IT或数字部门。报告通过相关性分析发现，CEO对AI治理的参与度，是所有25个测试属性中，与企业EBIT影响相关度最高的因素之一。在大公司样本中，这个相关性甚至更强。

KC评论： 这意味着AI不是“买了工具就完事”的IT项目，而是一场需要CEO亲自拉动的组织变革。报告没有明说但数据暗示的是：如果CEO不参与，AI项目大概率会变成“成本中心”而非“利润中心”。完整报告里有一张25个属性的相关性排序图，值得仔细看——它揭示了哪些投入真正能转化为利润。

报告测试了25个属性对企业EBIT影响的相关性，结果排名第一的是“是否重新设计了工作流”。不是买了多少GPU，不是用了哪个模型，而是企业是否真正改变了“人+AI”协同的方式。

目前只有21\%的受访者表示，他们的组织已经对至少部分工作流进行了根本性重新设计。这个数字本身就说明问题：大多数企业还是在“旧流程上贴AI标签”，而不是重新思考“这件事应该怎么做”。

报告引用了麦肯锡合伙人Alex Singla的评论：那些从AI中建立持久竞争优势的企业，思考的是“彻底的转型变革”，而不是“一个接一个的用例”。这种端到端的思维方式，决定了你是把AI当作修补工具，还是当作重塑商业模式的杠杆。

报告显示，28\%的受访者表示CEO负责AI治理，17\%表示由董事会负责。但更值得注意的是：平均而言，每个组织有两位领导共同负责AI治理。这不是“责任分散”，而是“风险共担”的信号。

在风险缓解方面，报告发现企业正在加速应对三类风险：不准确性、网络安全和知识产权侵权。与2024年初的调查相比，这三类风险的缓解措施都在显著增加。大型企业尤其积极，在网络安全和隐私风险方面投入明显更多。

但报告也揭示了一个值得警惕的差异：大型企业在处理“可解释性”和“准确性”风险上，并没有比小企业更积极。这暗示了一个潜在盲区——规模越大，AI输出的“黑箱问题”可能越被忽视。

KC评论： 这里有一个值得追问的伏笔：当AI输出的审查比例差异巨大时（27\%的企业审查全部输出，同样27\%的企业只审查20\%或更少），不同企业的“AI事故容忍度”可能完全不一样。完整报告里有一张按行业细分的审查比例图，能帮你判断自己所在行业处于哪个区间。

报告列出了12项AI采用与规模化的最佳实践，从“建立专门的AI推进团队”到“追踪明确定义的KPI”。相关性分析显示，每一项实践都与更高的EBIT影响正相关。

但现实是：不到三分之一的受访者表示他们的组织在遵循大多数实践。只有不到五分之一的企业在追踪AI项目的KPI。这意味着，大多数企业还在“凭感觉”管理AI投入。

最有趣的是，不同规模企业的差距：大型企业在“建立明确的规

[... middle omitted ...]

用率”“自动化率”还是“决策准确率”？不同行业、不同职能的KPI应该怎么差异化设计？

这可能是报告有意留下的“未完成讨论”。因为从麦肯锡的咨询实践来看，KPI设计恰恰是大多数企业最头疼的问题——设错了，AI项目就会变成“为了用AI而用AI”的表演。

另一个报告没有完全展开的问题是：当AI输出审查比例差异巨大时，不同企业的“风险-效率”平衡点在哪里？是“全部审查”更安全，还是“重点审查”更高效？报告给出了现状数据，但没有给出“最佳实践”建议。

KC评论： 这两个未解问题恰恰是这份报告最有价值的“后续阅读入口”。如果你正负责推动AI项目，完整报告里有12项实践的详细定义和按行业/规模的交叉分析，能直接帮你对照自己组织的成熟度。

基于报告数据，我们建议读者用以下三个维度来观察一家企业的AI成熟度，而不是只看“有没有用AI”：

*第一个维度：治理结构。** CEO是否亲自参与AI治理？董事会是否有AI监督机制？如果AI治理被下放到IT部门，大概率还处于“实验阶段”。

\section{Disclaimer}
{\scriptsize\color{refgray}
This community is a paid learning and information-sharing community created for general educational purposes only. The membership fee is charged solely for access to community discussions, educational content, organization, and operational costs. It is not an advisory fee, management fee, performance fee, brokerage fee, or compensation for personalized financial, investment, legal, tax, or professional advice.\par
I participate and share information solely as an individual and for educational discussion among community members. I am not acting as a financial adviser, investment adviser, broker, dealer, portfolio manager, legal adviser, tax adviser, accountant, fiduciary, or any other licensed professional adviser. Nothing shared in this community constitutes, and should not be interpreted as, financial advice, investment advice, legal advice, tax advice, a recommendation, solicitation, offer, endorsement, instruction, or guarantee to buy, sell, hold, short, trade, or invest in any security, cryptocurrency, fund, derivative, strategy, or other financial product.\par
All content shared in this community, including messages, discussions, links, files, charts, examples, market commentary, watchlists, AI-generated or AI-polished summaries, and third-party materials, is general, impersonal, and educational in nature. It does not take into account any individual's financial situation, investment objectives, risk tolerance, investment horizon, liquidity needs, tax status, legal restrictions, or suitability requirements. No content should be relied upon as suitable for any specific person, account, portfolio, or financial decision.\par
Information shared in this community may be incomplete, inaccurate, outdated, speculative, AI-generated, AI-edited, or based on third-party internet sources that have not been independently verified. I make no representation or warranty as to the accuracy, completeness, timeliness, reliability, or suitability of any information shared.\par
Financial markets involve substantial risk, including the possible loss of principal. Past performance, historical data, hypothetical examples, simulated results, backtests, personal opinions, or educational examples do not guarantee future results. Each member is solely responsible for conducting their own research, exercising independent judgment, and making their own decisions. Before making any financial, investment, legal, or tax decision, members should consult qualified and licensed professionals.\par
Participation in this community, payment of any membership fee, or communication with me or other members does not create any adviser-client, fiduciary, professional, contractual, agency, partnership, employment, or other advisory relationship. By joining or participating in this community, you acknowledge and agree that you are solely responsible for your own decisions, actions, transactions, profits, losses, risks, and consequences, and that I shall not be liable for any direct, indirect, incidental, consequential, financial, legal, tax, or other loss or damage arising from or related to any content, discussion, information, or material shared in this community.\par
}
\end{document}
